\documentclass[11pt,letterpaper]{article}

\usepackage[spanish,es-noshorthands]{babel}
\usepackage{fontspec}
\setmainfont{Source Serif Pro}
\usepackage{microtype}
\usepackage{amsmath}
\usepackage{geometry}
\usepackage{booktabs}
\usepackage{tabularx}
\usepackage{enumitem}
\usepackage{xcolor}
\usepackage{csquotes}
\usepackage{float}
\usepackage{tikz}
\usetikzlibrary{arrows.meta,positioning,shapes.geometric,calc,fit,backgrounds,shadows}
\usepackage{pgfplots}
\pgfplotsset{compat=1.18}
\usepackage{tocloft}
\usepackage{xurl}
\usepackage[hidelinks]{hyperref}

% Índice general compacto: el interlineado de \parskip lo haría desbordar
\renewcommand{\cfttoctitlefont}{\large\bfseries}
\setlength{\cftbeforesecskip}{0.4em}
\setlength{\cftbeforesubsecskip}{0.05em}
\setlength{\cftaftertoctitleskip}{0.9em}
\renewcommand{\cftsecaftersnum}{.}  % "9." como en los títulos de sección
\setlength{\cftsecnumwidth}{1.9em}
\usepackage[backend=biber,style=apa,sorting=nyt]{biblatex}
\DeclareLanguageMapping{spanish}{spanish-apa}

\geometry{margin=2.5cm}
\setlength{\emergencystretch}{1em}
\setlength{\parindent}{0pt}
\setlength{\parskip}{0.65em}
\setlist{nosep}
\newcommand{\pct}{\char`\%}  % \% sin el espacio fino de babel

% Apartado sin numerar, con entrada propia en el índice general
\newcommand{\subsec}[1]{%
  \subsection*{#1}%
  \phantomsection\addcontentsline{toc}{subsection}{#1}}
\addbibresource{referencias.bib}
\AtBeginBibliography{\sloppy}

% --------------------------------------------------------------------------
% Estilos gráficos comunes
% --------------------------------------------------------------------------
\definecolor{capa}{HTML}{2F5D8C}
\definecolor{trans}{HTML}{8C5A2F}
\definecolor{acento}{HTML}{9C2B2B}
\definecolor{jugA}{HTML}{1B7F9E}    % jugador A: Terraform
\definecolor{jugB}{HTML}{6F8F14}    % jugador B: Bicep
\definecolor{jugC}{HTML}{9C2B2B}    % jugador C: Pulumi
\definecolor{propia}{HTML}{2F5D8C}  % ruta propuesta para la empresa
\definecolor{descart}{HTML}{7A7A7A} % opción evaluada y descartada

\definecolor{opmobjc}{HTML}{1E7A1E}   % contorno de los objetos (OPM)
\definecolor{opmprocc}{HTML}{2E3FA8}  % contorno de los procesos (OPM)
\definecolor{dsmarq}{HTML}{2F5D8C}    % marcas de la DSM de arquitectura
\definecolor{dsmpro}{HTML}{8C5A2F}    % marcas de la DSM de cartera

\tikzset{
  % ---------------- Notación OPM (ISO 19450) ----------------
  % Objeto: rectángulo de contorno verde.
  opmobj/.style   = {draw=opmobjc, very thick, rectangle, fill=opmobjc!7,
                     font=\footnotesize\bfseries, align=center,
                     inner sep=4pt, minimum height=8mm, text width=26mm},
  opmobjn/.style  = {opmobj, text width=},
  % Proceso: elipse de contorno azul.
  opmproc/.style  = {draw=opmprocc, very thick, ellipse, fill=opmprocc!7,
                     font=\footnotesize\bfseries, align=center, text=black,
                     inner sep=1pt, minimum height=9mm, minimum width=28mm},
  % Estado: rectángulo de esquinas redondeadas, dentro del objeto.
  opmstate/.style = {draw=opmobjc, thick, rounded corners=3pt, fill=white,
                     font=\scriptsize, inner sep=2pt, minimum height=4.5mm},
  % Cosa física: sombra proyectada. Cosa ambiental: contorno discontinuo.
  opmfis/.style   = {drop shadow={shadow xshift=0.8mm, shadow yshift=-0.8mm,
                     opacity=1, fill=black!45}},
  opmamb/.style   = {densely dashed},
  % Enlaces procedimentales.
  instr/.style    = {thick, -{Circle[open,length=5pt]}},   % instrumento
  agente/.style   = {thick, -{Circle[length=5pt]}},        % agente
  efecto/.style   = {thick, {Stealth[length=5.5pt]}-{Stealth[length=5.5pt]}},
  flecha/.style   = {thick, -{Stealth[length=5pt]}},       % consumo/resultado
  invoca/.style   = {thick, densely dashed, -{Stealth[length=5pt]}},
  % Diagramas de bloques
  bloque/.style   = {draw, thick, rounded corners=2pt, align=center,
                     font=\footnotesize, inner sep=5pt},
  etiq/.style     = {font=\scriptsize\itshape, align=center},
}

% --------------------------------------------------------------------------
% Símbolos estructurales de OPM.
% La base del triángulo se apoya en el lado del todo, del general o del
% exhibidor, y el vértice apunta hacia la parte, la especialización o el
% atributo. #1 = rotación en grados (0 = base arriba, vértice abajo);
% #2 = coordenada previamente definida sobre la que se centra el símbolo.
% --------------------------------------------------------------------------
\newcommand{\opmtri}[4]{% #3 = semiancho, #4 = relleno
  \begin{scope}[rotate around={#1:(#2)}]
    \path[draw,very thick,fill=#4] ([shift={(-#3,#3)}]#2) --
      ([shift={(#3,#3)}]#2) -- ([shift={(0mm,-#3)}]#2) -- cycle;
  \end{scope}}
\newcommand{\opmagg}[2]{\opmtri{#1}{#2}{2.1mm}{black}}   % agregación-participación
\newcommand{\opmgen}[2]{\opmtri{#1}{#2}{2.1mm}{white}}   % generalización-especialización
\newcommand{\opmexh}[2]{% exhibición-caracterización: triángulo dentro de triángulo
  \opmtri{#1}{#2}{2.6mm}{white}%
  \begin{scope}[rotate around={#1:(#2)}]
    \path[fill=black] ([shift={(-1.1mm,0.9mm)}]#2) --
      ([shift={(1.1mm,0.9mm)}]#2) -- ([shift={(0mm,-1.1mm)}]#2) -- cycle;
  \end{scope}}

% --------------------------------------------------------------------------
% Matriz de estructura de dependencia (DSM).
% Rejilla cuadrada de lado #1 celdas, con celda de \dsmu de lado.
% --------------------------------------------------------------------------
\newlength{\dsmu}\setlength{\dsmu}{6.4mm}

\title{\textbf{Alcance de la tecnología: la infraestructura como código}\\
\large Funcionamiento de un sistema completo, capas y módulos principales,
matriz de estructura de dependencia, modelo de objetos y procesos, y
posicionamiento frente a la competencia}
\author{Carlos Luis Rojas Aragonés\\
Gestión del Desarrollo Tecnológico}
\date{1 de septiembre de 2026}

\begin{document}

\maketitle

\begin{abstract}
\noindent\small
La infraestructura como código (\emph{Infrastructure as Code}, IaC) sustituye
la configuración manual de servidores, redes y servicios por una
especificación versionada que un motor automático traduce en operaciones sobre
la interfaz de programación de un proveedor. Este trabajo delimita el alcance
de la tecnología, describe su principio de funcionamiento (la conciliación
continua entre un estado deseado y un estado real), desglosa las siete capas y
los seis módulos transversales de un sistema completo, y ordena sus
interdependencias en dos matrices de estructura de dependencia, una de
arquitectura y otra de cartera de proyectos, cuyo particionado produce la
secuencia de implantación. Formaliza después el conjunto mediante un modelo de
objetos y procesos según la metodología OPM, con su diagrama de sistema, el
despliegue estructural de la especificación, la ampliación en detalle del
proceso principal y el lenguaje OPL de los tres diagramas. Sobre las figuras
de mérito de la tecnología se construye por último el posicionamiento de la
empresa frente a la competencia, con un
gráfico vectorial que sitúa a Terraform, Bicep y Pulumi en un plano de
variación de valor para el cliente y variación de coste para el productor, y
que descompone la trayectoria de cada jugador en los vectores tecnológicos que
la producen. Todas las figuras son de elaboración propia.
\end{abstract}

\clearpage
\setcounter{tocdepth}{2}
{\setlength{\parskip}{0pt}\tableofcontents}
\clearpage

% ==========================================================================
\section{Objeto y delimitación del alcance}
\label{sec:alcance}
% ==========================================================================

\subsec{01. El encargo}

El documento recoge entregas encadenadas de una misma hoja de ruta
tecnológica, y por eso no se organiza por actividades sino por los elementos
del índice recomendado para una hoja de ruta. El Cuadro~\ref{tab:atra} recoge
esa correspondencia.

La primera entrega, el alcance de la tecnología, cubre los tres primeros
elementos. El elemento 1, la descripción de la tecnología, ocupa las
secciones~\ref{sec:principio} a \ref{sec:modulo}, con el resumen funcional y
los diagramas de capas y de módulos. El elemento 2, la matriz de estructura de
dependencia, ocupa la sección~\ref{sec:dsm} con dos matrices, una de los
componentes del sistema y otra de los proyectos de la cartera, y con el
particionado que convierte a ambas en una secuencia de implantación. El
elemento 3, el modelo de objetos y procesos, ocupa la
sección~\ref{sec:opm}, con el diagrama de sistema, el despliegue estructural
de la especificación, la ampliación en detalle del proceso principal y el OPL
de los tres diagramas.

La segunda entrega pide analizar la competencia mediante un gráfico vectorial
que incluya al menos dos competidores. Corresponde al elemento 6 y ocupa la
sección~\ref{sec:comp}, que compara tres jugadores, Terraform, Bicep y
Pulumi, con la trayectoria propuesta para la empresa. Su sitio natural es
inmediatamente después de la sección~\ref{sec:fom}, porque las coordenadas del
gráfico se construyen a partir de las figuras de mérito definidas allí.

\begin{table}[ht]
  \centering
  \footnotesize
  \caption{Correspondencia entre los doce elementos del índice recomendado
  para una hoja de ruta tecnológica y las secciones de este documento.}
  \label{tab:atra}
  \begin{tabularx}{\textwidth}{@{}l>{\raggedright\arraybackslash}X
    >{\raggedright\arraybackslash}p{0.30\textwidth}@{}}
    \toprule
    \textbf{N.\textsuperscript{o}} & \textbf{Elemento} &
      \textbf{Dónde se atiende} \\
    \midrule
    1 & Descripción de la tecnología &
      Secciones \ref{sec:alcance} a \ref{sec:modulo} \\
    2 & Matriz de estructura de dependencia (DSM) &
      Sección \ref{sec:dsm} \\
    3 & Modelo con metodología OPM &
      Sección \ref{sec:opm} \\
    4 & Figuras de mérito (FOM) &
      Sección \ref{sec:fom} \\
    5 & Correspondencia con las motivaciones estratégicas &
      Pendiente \\
    6 & Posicionamiento frente a la competencia &
      Sección \ref{sec:comp} \\
    7 & Modelo técnico & Pendiente \\
    8 & Modelo financiero &
      Declarado fuera de alcance en la sección \ref{sec:comp} \\
    9 & Portfolio de proyectos y prototipos de I+D &
      Cuadro \ref{tab:dsmproy}, en lo relativo a su secuencia \\
    10 & Publicaciones, presentaciones y patentes clave & Pendiente \\
    11 & Declaración de estrategia tecnológica & Pendiente \\
    12 & Evaluación del nivel de madurez de la hoja de ruta & Pendiente \\
    \bottomrule
  \end{tabularx}
\end{table}

\subsec{02. Qué se entiende aquí por sistema completo}

Conviene fijar el término antes de dibujar nada, porque \enquote{infraestructura
como código} se usa con dos alcances muy distintos. En el sentido estrecho
designa una herramienta: un fichero declarativo y un binario que lo aplica. En
el sentido amplio, que es el que aquí interesa, designa un sistema
sociotécnico completo, con su cadena de herramientas, su modelo de estado, sus
controles de gobierno y el equipo que lo opera.

Un sistema completo de IaC es aquel que satisface cuatro propiedades
simultáneas \parencite{morris2021}:

\begin{enumerate}
  \item \textbf{Especificación externalizada.} La configuración deseada de la
    infraestructura vive fuera de la infraestructura, en artefactos de texto
    legibles y versionados, no en la memoria de un operador ni en el estado
    acumulado de las máquinas.
  \item \textbf{Reproducibilidad.} Aplicar la misma especificación sobre un
    sustrato vacío produce un resultado equivalente, con independencia de
    quién la aplique y de cuándo lo haga.
  \item \textbf{Idempotencia.} Aplicar la especificación sobre un sistema que
    ya la cumple no produce ningún cambio, y aplicarla \emph{n} veces tiene el
    mismo efecto que aplicarla una vez.
  \item \textbf{Conciliación.} El sistema es capaz de detectar la divergencia
    entre lo especificado y lo realmente desplegado, y de corregirla.
\end{enumerate}

La literatura académica sobre la materia, todavía escasa en comparación con su
difusión industrial, se ha concentrado sobre todo en la calidad y los defectos
del propio código de infraestructura \parencite{rahman2019}.

Un montaje que cumpla sólo la primera propiedad es un repositorio de guiones,
no un sistema de IaC. La cuarta es la que convierte la automatización en
control, y es también la que con más frecuencia falta en las implantaciones
reales.

\subsec{03. La frontera del sistema}

Delimitar el alcance exige decir también qué queda fuera. El
Cuadro~\ref{tab:frontera} fija la frontera adoptada en este documento.

\begin{table}[ht]
  \centering
  \footnotesize
  \caption{Frontera del sistema de infraestructura como código.}
  \label{tab:frontera}
  \begin{tabularx}{\textwidth}{@{}>{\raggedright\arraybackslash}X
    >{\raggedright\arraybackslash}X@{}}
    \toprule
    \textbf{Dentro del alcance} & \textbf{Fuera del alcance} \\
    \midrule
    Cómputo, red, almacenamiento, identidad y servicios gestionados
      declarados como código. &
    El sustrato físico del proveedor: centros de datos, cableado,
      refrigeración. \\
    Configuración del sistema operativo y de los agentes instalados en él. &
    El código de la aplicación de negocio y sus pruebas funcionales. \\
    Construcción de imágenes y artefactos inmutables. &
    Los datos de producción y su ciclo de vida, que exigen migraciones y
      copias de seguridad con reglas propias. \\
    Orquestación de cargas y despliegue continuo de manifiestos. &
    La gestión de incidencias, la guardia y los procedimientos de crisis. \\
    Estado, secretos, políticas, pruebas y observabilidad de la propia
      infraestructura. &
    La observabilidad funcional de la aplicación (trazas de negocio,
      experiencia de usuario). \\
    \bottomrule
  \end{tabularx}
\end{table}

La distinción entre infraestructura y datos merece un comentario, porque es la
frontera que más problemas causa en la práctica. Un servidor de base de datos
está dentro del alcance: su tamaño, su versión, su red y sus copias
programadas se declaran como código. El contenido de esa base de datos no lo
está: no es idempotente, no se puede destruir y recrear, y su presencia obliga
a que el sistema de IaC distinga recursos desechables de recursos con estado
persistente, y proteja a los segundos de las operaciones de reemplazo.

% ==========================================================================
\section{El principio de funcionamiento}
\label{sec:principio}
% ==========================================================================

\subsec{01. Tres estados y una función}

Todo el funcionamiento de un sistema de IaC se sostiene sobre tres
representaciones de la misma infraestructura y sobre la función que las
compara.

\begin{description}[style=nextline, leftmargin=0pt, labelwidth=0pt, labelsep=0pt]
  \item[Estado deseado.] Lo que la especificación declara. Es un artefacto de
    texto, versionado, que describe recursos y sus propiedades. No describe
    los pasos para llegar a ellos.
  \item[Estado registrado.] Lo que el motor cree que existe, resultado de la
    última operación que él mismo ejecutó. Se materializa en un fichero o
    servicio de estado que asocia cada recurso declarado con el identificador
    real que el proveedor le asignó \parencite{terraformstate}.
  \item[Estado real.] Lo que efectivamente existe en el proveedor, que sólo se
    conoce consultando su interfaz de programación.
\end{description}

La función central del sistema es la \textbf{conciliación}: calcular la
diferencia entre el estado deseado y el estado real, y emitir el conjunto
mínimo de operaciones que anula esa diferencia. La Figura~\ref{fig:bucle}
representa el bucle.

\begin{figure}[ht]
  \centering
  \begin{tikzpicture}
    \node[bloque, fill=blue!8, text width=30mm, minimum height=12mm] (des)
      {\textbf{Estado deseado}\\\scriptsize especificación versionada};
    \node[bloque, fill=orange!12, text width=30mm, minimum height=12mm,
      right=26mm of des] (comp)
      {\textbf{Comparador}\\\scriptsize cálculo de la diferencia};
    \node[bloque, fill=green!10, text width=30mm, minimum height=12mm,
      right=26mm of comp] (real)
      {\textbf{Estado real}\\\scriptsize recursos del proveedor};
    \node[bloque, fill=gray!10, text width=30mm, minimum height=12mm,
      below=17mm of comp] (act)
      {\textbf{Actuador}\\\scriptsize crear, modificar, destruir};

    \draw[flecha] (des) -- node[above, etiq] {declaración} (comp);
    \draw[flecha] (real) -- node[above, etiq] {observación} (comp);
    \draw[flecha] (comp) -- node[right, etiq, xshift=1mm] {plan} (act);
    \draw[flecha] (act.east) to[out=0, in=-90]
      node[pos=0.55, below right, etiq, xshift=-3mm, yshift=1mm]
      {\parbox{24mm}{\centering operaciones\\sobre la API}} (real.south);
  \end{tikzpicture}
  \caption{Bucle de conciliación. El sistema no ejecuta una secuencia de
  pasos, sino que cierra repetidamente la distancia entre dos
  representaciones. Si la diferencia es vacía, el actuador no ejecuta nada, y
  ésa es exactamente la propiedad de idempotencia. Elaboración propia.}
  \label{fig:bucle}
\end{figure}

\subsec{02. Declarativo frente a imperativo}

La diferencia entre los dos paradigmas no es de sintaxis sino de quién calcula
el camino. En el enfoque imperativo, el autor escribe la secuencia de
operaciones y asume la responsabilidad de que sea correcta desde cualquier
estado de partida. En el enfoque declarativo, el autor escribe el destino y el
motor deriva la secuencia.

La consecuencia práctica es que un guion imperativo sólo es correcto para el
estado de partida que su autor imaginó, mientras que una especificación
declarativa es correcta para cualquiera, a condición de que el motor sepa
observar el estado real. Por eso el paradigma declarativo domina la capa de
aprovisionamiento, donde el espacio de estados de partida es enorme, y
sobrevive el imperativo en tareas de configuración muy locales, donde ese
espacio es pequeño y conocido.

Casi ningún sistema es puro. Las herramientas de configuración expresan los
pasos en orden, pero cada paso es un módulo idempotente que sólo actúa si hace
falta \parencite{ansibledocs}. Es un híbrido: secuencia imperativa de
operaciones declarativas.

\subsec{03. Convergencia y deriva}

Un sistema convergente es el que, aplicado repetidamente, se aproxima al
estado deseado y luego se queda quieto. La \textbf{deriva}
(\emph{drift}) es la aparición de diferencias entre el estado real y el
deseado por causas ajenas al sistema: un cambio manual de urgencia, una acción
automática del proveedor, la caducidad de un certificado, el borrado
accidental de un recurso.

La deriva es inevitable, y la madurez de una implantación se mide por lo que
hace con ella. Hay tres posturas, en orden creciente de rigor:

\begin{enumerate}
  \item \textbf{Ignorarla.} El motor sólo se ejecuta cuando alguien lo lanza.
    Entre ejecuciones, el estado real puede alejarse arbitrariamente.
  \item \textbf{Detectarla.} Una ejecución periódica en modo de sólo lectura
    compara y notifica, sin corregir. Convierte la deriva en una métrica
    observable.
  \item \textbf{Corregirla.} Un agente concilia de forma continua y revierte
    todo cambio no declarado. Es el principio de conciliación continua de
    GitOps \parencite{opengitops} y el modo nativo de los controladores de
    Kubernetes \parencite{k8scontrollers,burns2016}.
\end{enumerate}

% ==========================================================================
\section{El ciclo básico de operación}
\label{sec:ciclo}
% ==========================================================================

Descendiendo un nivel desde el bucle abstracto, la operación cotidiana de un
sistema completo recorre siete pasos. La Figura~\ref{fig:ciclo} los ordena.

\begin{figure}[ht]
  \centering
  \begin{tikzpicture}[node distance=4.5mm and 5.5mm]
    \tikzset{paso/.style={bloque, fill=capa!12, draw=capa, text width=20mm,
      minimum height=13mm, font=\scriptsize}}
    \node[paso] (p1) {\textbf{1. Codificar}\\ escribir la\\especificación};
    \node[paso, right=of p1] (p2) {\textbf{2. Validar}\\ formato, tipos,\\
      políticas};
    \node[paso, right=of p2] (p3) {\textbf{3. Planificar}\\ diferencia\\
      deseado--real};
    \node[paso, right=of p3] (p4) {\textbf{4. Autorizar}\\ revisión y\\
      aprobación};
    \node[paso, right=of p4] (p5) {\textbf{5. Aplicar}\\ ejecutar sobre\\
      la API};
    \node[paso, below=9mm of p4] (p6) {\textbf{6. Verificar}\\ pruebas de\\
      infraestructura};
    \node[paso, below=9mm of p2] (p7) {\textbf{7. Conciliar}\\ vigilar y\\
      corregir la deriva};

    \foreach \a/\b in {p1/p2,p2/p3,p3/p4,p4/p5}
      \draw[flecha, draw=capa] (\a) -- (\b);
    \draw[flecha, draw=capa] (p5.south) |- (p6.east);
    \draw[flecha, draw=capa] (p6) -- (p7);
    \draw[flecha, draw=acento, dashed] (p7.west) -|
      node[pos=0.28, below, etiq, text=acento, yshift=-1mm]
      {deriva detectada} (p1.south);
  \end{tikzpicture}
  \caption{Los siete pasos del ciclo de operación. El retorno discontinuo es
  el que distingue un sistema de IaC de una simple automatización de
  despliegue. Elaboración propia.}
  \label{fig:ciclo}
\end{figure}

Los pasos 2, 4 y 6 son los que suelen faltar en implantaciones incompletas, y
son precisamente los que aportan garantía. Sin validación de políticas, el
sistema automatiza también los errores de configuración; sin autorización,
elimina el punto de control humano sobre cambios irreversibles; sin
verificación, no distingue entre \enquote{el motor terminó sin error} y
\enquote{la infraestructura funciona}.

\subsec{01. El paso crítico: la comparación a tres bandas}

El paso 3 merece detalle porque en él se concentra el mecanismo que hace útil
al sistema. El motor no compara dos cosas sino tres, y lo hace en dos fases
(Figura~\ref{fig:tresbandas}).

\begin{figure}[ht]
  \centering
  \begin{tikzpicture}
    \tikzset{caja/.style={bloque, text width=34mm, minimum height=15mm}}
    \node[caja, fill=blue!8]  (d) {\textbf{Deseado}\\\scriptsize lo que
      declara el código};
    \node[caja, fill=gray!12, right=20mm of d] (r) {\textbf{Registrado}\\
      \scriptsize lo que el motor anotó\\ en la última ejecución};
    \node[caja, fill=green!10, right=20mm of r] (a) {\textbf{Real}\\
      \scriptsize lo que devuelve\\ la API del proveedor};

    \draw[efecto, draw=trans] (r) -- node[above, etiq, text=trans]
      {\textbf{fase 1}} node[below, etiq, text=trans] {refresco} (a);
    \draw[efecto, draw=acento] (d) -- node[above, etiq, text=acento]
      {\textbf{fase 2}} node[below, etiq, text=acento] {plan} (r);

    \node[bloque, fill=acento!10, draw=acento, below=12mm of r,
      text width=0.62\textwidth]
      {\textbf{Plan de cambios}: lista de recursos a crear, modificar,
       reemplazar o destruir, con el valor anterior y el posterior de cada
       propiedad afectada};
    \draw[flecha, draw=acento] ($(d.south)!0.5!(a.south)+(0,-2mm)$) --
      ++(0,-5mm);
  \end{tikzpicture}
  \caption{Comparación a tres bandas. La fase de refresco actualiza el
  registro contra el mundo real y es la que descubre la deriva; la fase de
  plan compara el código contra ese registro ya actualizado y es la que
  produce las operaciones. Elaboración propia.}
  \label{fig:tresbandas}
\end{figure}

El estado registrado no es redundante. Cumple tres funciones que ninguna de
las otras dos representaciones puede cumplir. Primera, traduce el nombre
lógico que el código da a un recurso al identificador opaco que el proveedor
le asignó, sin lo cual el motor no podría reconocer que un recurso existente
le pertenece. Segunda, permite detectar destrucciones: un recurso presente en
el registro y ausente del proveedor sólo puede detectarse comparando contra el
registro, nunca comparando el código con la realidad. Tercera, guarda
propiedades calculadas por el proveedor que el código no declara y que otros
recursos necesitan.

Ese fichero es, en consecuencia, un activo crítico: contiene identificadores y
con frecuencia valores sensibles, y su corrupción o pérdida deja al sistema
sin capacidad de reconocer lo que él mismo creó. De ahí que todo montaje serio
lo almacene en un servicio remoto con cifrado, versionado y bloqueo de
concurrencia.

% ==========================================================================
\section{Las capas de un sistema completo}
\label{sec:capas}
% ==========================================================================

Un sistema completo se organiza en siete capas apiladas, cada una de las
cuales consume el resultado de la anterior y expone una abstracción más alta.
La Figura~\ref{fig:pila} las presenta con sus herramientas representativas.

\begin{figure}[ht]
  \centering
  \begin{tikzpicture}[y=13.6mm]
    \def\W{14.6}
    \foreach \i/\nom/\fun/\tool in {
      0/{L0 · Sustrato}/{cómputo, red, almacenamiento e identidad del
         proveedor}/{API de AWS, Azure o GCP; equipo físico},
      1/{L1 · Aprovisionamiento}/{crea, modifica y destruye los recursos
         declarados}/{Terraform, OpenTofu, Pulumi, CloudFormation, Bicep},
      2/{L2 · Imágenes}/{construye artefactos inmutables ya
         configurados}/{Packer, Dockerfile, \emph{buildpacks}},
      3/{L3 · Configuración}/{ajusta el interior de la máquina ya
         arrancada}/{Ansible, Chef, Puppet, cloud-init},
      4/{L4 · Orquestación}/{programa cargas sobre el conjunto de
         máquinas}/{Kubernetes, Nomad, Helm, Kustomize},
      5/{L5 · Entrega}/{lleva el código al motor y concilia de forma
         continua}/{Argo CD, Flux, canalizaciones de integración},
      6/{L6 · Plataforma}/{expone catálogo y autoservicio al equipo de
         producto}/{Backstage, Crossplane, portales internos}} {
      \node[draw=capa, very thick, fill=capa!\the\numexpr 6+2*\i\relax,
        rounded corners=2pt, minimum width=\W cm, minimum height=12.6mm,
        anchor=south west] (L\i) at (0,\i) {};
      \node[anchor=west, font=\footnotesize\bfseries, text=capa!60!black,
        text width=37mm, align=left] at ([xshift=3.5mm]L\i.west) {\nom};
      \node[anchor=west, font=\scriptsize, text width=50mm, align=left]
        at ([xshift=43mm]L\i.west) {\fun};
      \node[anchor=west, font=\scriptsize\itshape, text width=45mm,
        align=left] at ([xshift=96mm]L\i.west) {\tool};
    }
    \draw[flecha, draw=capa, line width=1pt] (-0.5,0.2) -- (-0.5,6.8)
      node[midway, rotate=90, above, font=\scriptsize, text=capa]
      {abstracción creciente};
  \end{tikzpicture}
  \caption{Las siete capas de un sistema completo de infraestructura como
  código. Cada capa consume la salida de la inferior. Las herramientas citadas
  son representativas, no exhaustivas. Elaboración propia.}
  \label{fig:pila}
\end{figure}

\subsec{01. Qué hace cada capa}

\textbf{L0, sustrato.} No es parte del sistema de IaC sino su objeto. Aporta
la interfaz de programación sobre la que todo lo demás actúa. Su
característica determinante es que ofrece recursos con identidad propia,
creables y destruibles por llamada, lo que es la condición de posibilidad de
todo el resto.

\textbf{L1, aprovisionamiento.} Traduce la especificación en llamadas de
creación, modificación y destrucción. Es la capa que posee el registro de
estado y la que ejecuta la comparación a tres bandas. Trabaja con recursos que
todavía no tienen sistema operativo arrancado, o cuyo interior no le
concierne. Conviene señalar que en 2023 el cambio de licencia de Terraform dio
origen a OpenTofu, una bifurcación acogida por la Linux Foundation
\parencite{opentofu}, de modo que hoy la capa cuenta con dos implementaciones
compatibles del mismo lenguaje: una decisión de dependencia que una hoja de
ruta tecnológica debería registrar de forma explícita.

\textbf{L2, imágenes.} Desplaza el trabajo de configuración desde el momento
de ejecución al momento de construcción. En lugar de arrancar una máquina
genérica y modificarla, construye una imagen ya completa y arranca instancias
idénticas de ella. Es la base de la infraestructura inmutable: en lugar de
modificar un servidor, se reemplaza. Reduce drásticamente la deriva, porque
elimina la ventana temporal en la que un servidor puede divergir.

\textbf{L3, configuración.} Actúa sobre el interior de máquinas ya arrancadas:
paquetes, ficheros, servicios, usuarios. Es la capa históricamente más antigua
y la que más terreno ha cedido a L2 y L4. Sigue siendo necesaria donde hay
máquinas de larga vida, sistemas heredados o equipos físicos.

\textbf{L4, orquestación.} Introduce una indirección importante: deja de
declararse \enquote{esta carga en esta máquina} y pasa a declararse
\enquote{esta carga, tantas réplicas} sobre un conjunto de máquinas
intercambiables. El planificador decide la colocación. Sus controladores son
bucles de conciliación como el de la Figura~\ref{fig:bucle}, pero embebidos y
permanentes \parencite{k8scontrollers}.

\textbf{L5, entrega.} Conecta el repositorio con el motor. Ejecuta las
canalizaciones que validan, planifican y aplican, y, en el modelo de
conciliación continua, aloja el agente que vigila la deriva sin intervención
humana \parencite{opengitops,humble2010}.

\textbf{L6, plataforma.} Envuelve todo lo anterior en una interfaz de
autoservicio, para que un equipo de producto pueda solicitar una base de datos
sin conocer ni la capa L1 ni las políticas de la organización, que quedan
codificadas en el catálogo.

\subsec{02. La tensión entre capas}

Las capas L2, L3 y L4 compiten parcialmente por el mismo territorio. Un mismo
requisito, por ejemplo \enquote{que este servidor tenga la versión 1.24 de un
tiempo de ejecución}, puede resolverse construyendo una imagen que ya la
incluya (L2), instalándola con un módulo de configuración (L3) o empaquetando
la aplicación en un contenedor que la lleve dentro (L4). Las tres soluciones
funcionan y no deben coexistir sobre el mismo recurso: si dos capas gobiernan
la misma propiedad, cada una revierte el trabajo de la otra en cada ciclo de
conciliación, y el sistema oscila.

Decidir qué capa posee cada propiedad es, por tanto, la primera decisión de
arquitectura de una implantación, y conviene documentarla de forma explícita.

% ==========================================================================
\section{Los módulos transversales}
\label{sec:transversales}
% ==========================================================================

Las siete capas describen el flujo principal, pero un sistema completo
necesita además seis módulos que atraviesan todas las capas
(Figura~\ref{fig:transversales}). Ninguno produce infraestructura; todos
condicionan si la que se produce es correcta, segura y auditable.

\begin{figure}[ht]
  \centering
  \begin{tikzpicture}
    \node[draw=capa, very thick, fill=capa!10, rounded corners=3pt,
      minimum width=5.4cm, minimum height=3.7cm, align=center,
      font=\footnotesize\bfseries, text=capa!60!black] (nucleo)
      {Las siete capas\\del flujo principal\\[2pt]
       \scriptsize\mdseries L0 sustrato · L1 aprovisionamiento\\
       \scriptsize\mdseries L2 imágenes · L3 configuración\\
       \scriptsize\mdseries L4 orquestación · L5 entrega\\
       \scriptsize\mdseries L6 plataforma};
    \foreach \ang/\nom/\det in {
      90/{Control de versiones}/{historia, revisión, reversión},
      32/{Gestión de estado}/{registro, bloqueo, cifrado},
      -32/{Gestión de secretos}/{inyección en tiempo de ejecución},
      -90/{Política como código}/{reglas evaluables antes de aplicar},
      -148/{Pruebas}/{estáticas, unitarias, de integración},
      148/{Observabilidad}/{deriva, coste, cumplimiento}} {
      \node[bloque, fill=trans!12, draw=trans, text width=31mm,
        font=\scriptsize, align=center, anchor=center] at (\ang:5.1cm)
        (m\ang) {\textbf{\nom}\\ \det};
      \draw[efecto, draw=trans, shorten >=3pt, shorten <=3pt]
        (m\ang) -- (nucleo);
    }
  \end{tikzpicture}
  \caption{Los seis módulos transversales. Cada uno interactúa con todas las
  capas del flujo principal, no con una en particular. Elaboración propia.}
  \label{fig:transversales}
\end{figure}

El Cuadro~\ref{tab:transversales} resume la función de cada uno y el fallo
característico que aparece cuando falta.

\begin{table}[ht]
  \centering
  \footnotesize
  \caption{Módulos transversales: función, instrumentos y fallo característico
  por ausencia.}
  \label{tab:transversales}
  \begin{tabularx}{\textwidth}{@{}>{\raggedright\arraybackslash}p{0.145\textwidth}
    >{\raggedright\arraybackslash}X
    >{\raggedright\arraybackslash}p{0.30\textwidth}@{}}
    \toprule
    \textbf{Módulo} & \textbf{Función} & \textbf{Fallo por ausencia} \\
    \midrule
    Control de versiones & Registrar quién cambió qué, cuándo y por qué;
      permitir revertir a un estado anterior conocido. &
      Cambios sin trazabilidad y reversión imposible. \\
    Gestión de estado & Custodiar el registro: almacenamiento remoto, cifrado,
      versionado y bloqueo para impedir ejecuciones concurrentes. &
      Corrupción del registro y recursos huérfanos. \\
    Gestión de secretos & Mantener credenciales fuera del código e inyectarlas
      en el momento de la ejecución, con rotación y ámbito mínimo. Es uno de
      los defectos más frecuentes en código de infraestructura
      \parencite{rahman2019smells}. &
      Credenciales escritas en el repositorio. \\
    Política como código & Expresar las reglas de la organización como
      predicados evaluables sobre el plan, antes de aplicarlo
      \parencite{opa}. &
      El sistema automatiza también los incumplimientos, y a mayor
      velocidad. \\
    Pruebas & Verificar la especificación en varios niveles: análisis
      estático, validación del plan, despliegue efímero y comprobación
      funcional. &
      \enquote{Terminó sin error} se confunde con \enquote{funciona}. \\
    Observabilidad & Medir deriva, coste, tiempo de ciclo y cumplimiento, y
      exponerlos como series temporales. &
      No hay figuras de mérito y la mejora no es demostrable. \\
    \bottomrule
  \end{tabularx}
\end{table}

% ==========================================================================
\section{Modos de operación}
\label{sec:modos}
% ==========================================================================

Sobre la misma arquitectura caben decisiones de operación que cambian de forma
apreciable el comportamiento del sistema. Las tres principales son
independientes entre sí.

\subsec{01. Empuje frente a arrastre}

En el modo de \textbf{empuje} (\emph{push}), un actor externo, típicamente una
canalización de integración, posee las credenciales del entorno y ejecuta el
motor contra él. En el modo de \textbf{arrastre} (\emph{pull}), un agente
residente dentro del entorno consulta el repositorio y se aplica a sí mismo
los cambios. La Figura~\ref{fig:pushpull} contrasta ambos.

\begin{figure}[ht]
  \centering
  \begin{tikzpicture}
    \tikzset{c/.style={bloque, text width=28mm, minimum height=9mm},
             cin/.style={bloque, fill=white, text width=30mm,
               minimum height=7mm, font=\scriptsize},
             ent/.style={draw, thick, rounded corners=2pt, fill=green!10,
               minimum width=52mm}}
    % ---------------- empuje ----------------
    \node[etiq, font=\footnotesize\bfseries] (t1) {Empuje (\emph{push})};
    \node[c, fill=blue!8, below=5mm of t1] (r1) {Repositorio};
    \node[c, fill=orange!12, below=11mm of r1] (ci1)
      {Canalización\\ \scriptsize con las credenciales};
    \node[ent, minimum height=24mm, below=11mm of ci1] (e1) {};
    \node[anchor=north, font=\scriptsize\bfseries]
      at ([yshift=-2mm]e1.north) {Entorno};
    \node[cin, anchor=south] at ([yshift=3mm]e1.south) (i1)
      {Infraestructura};
    \draw[flecha] (r1) -- node[right, etiq, xshift=1mm] {disparo} (ci1);
    \draw[flecha] (ci1) -- node[right, etiq, xshift=1mm] {aplica} (e1.north);
    % ---------------- arrastre ----------------
    \node[etiq, font=\footnotesize\bfseries, right=52mm of t1] (t2)
      {Arrastre (\emph{pull})};
    \node[c, fill=blue!8, below=5mm of t2] (r2) {Repositorio};
    \node[ent, minimum height=34mm, below=26mm of r2] (e2) {};
    \node[anchor=north, font=\scriptsize\bfseries]
      at ([yshift=-1.5mm]e2.north) {Entorno};
    \node[cin, fill=orange!12, anchor=north] at ([yshift=-7mm]e2.north) (ag)
      {Agente\\ \scriptsize con las credenciales};
    \node[cin, anchor=south] at ([yshift=3mm]e2.south) (i2)
      {Infraestructura};
    \draw[flecha] ([xshift=-11mm]e2.north) to[out=95,in=-95]
      node[left, etiq, xshift=-1mm] {consulta} ([xshift=-9mm]r2.south);
    \draw[flecha] ([xshift=9mm]r2.south) to[out=-85,in=85]
      node[right, etiq, xshift=1mm] {\parbox{18mm}{\centering estado\\
      deseado}} ([xshift=11mm]e2.north);
    \draw[flecha, draw=acento] (ag) -- node[right, font=\tiny\itshape,
      text=acento, xshift=0.5mm] {concilia} (i2);
  \end{tikzpicture}
  \caption{Empuje frente a arrastre. En el arrastre, las credenciales del
  entorno nunca salen de él, lo que reduce la superficie de ataque y hace del
  agente el punto natural de conciliación continua. Elaboración propia a
  partir de los principios de \textcite{opengitops}.}
  \label{fig:pushpull}
\end{figure}

El arrastre tiene dos ventajas estructurales. La primera es de seguridad: las
credenciales de producción no viajan a un sistema de integración compartido.
La segunda es que el agente, por estar siempre presente, puede conciliar de
forma continua y no sólo cuando alguien confirma un cambio, que es la cuarta
propiedad exigida en la sección~\ref{sec:alcance}. Su coste es que exige un agente por
entorno y complica los flujos que necesitan orden entre entornos.

\subsec{02. Mutable frente a inmutable}

En el modelo mutable, un recurso desplegado se modifica en su sitio. En el
inmutable, no se modifica nunca: se construye uno nuevo con la configuración
deseada y se sustituye. El modelo inmutable elimina por construcción la deriva
de configuración y hace que el recurso desplegado sea siempre trazable a un
artefacto concreto. A cambio, obliga a que todo estado persistente viva fuera
del recurso reemplazable, lo que es una restricción de diseño exigente.

\subsec{03. El cuadro de decisión}

Las tres decisiones son independientes entre sí, de modo que caben ocho
combinaciones. La más frecuente en implantaciones recientes es la de arrastre,
inmutable y declarativa, pero ninguna de las tres opciones es obligada. El
Cuadro~\ref{tab:modos} resume las consecuencias de cada una.

\begin{table}[!ht]
  \centering
  \footnotesize
  \caption{Tres decisiones de operación y sus consecuencias.}
  \label{tab:modos}
  \begin{tabularx}{\textwidth}{@{}>{\raggedright\arraybackslash}p{0.15\textwidth}
    >{\raggedright\arraybackslash}X>{\raggedright\arraybackslash}X@{}}
    \toprule
    \textbf{Decisión} & \textbf{Opción A} & \textbf{Opción B} \\
    \midrule
    Iniciativa &
      \textbf{Empuje}: control de orden entre entornos; credenciales fuera del
      entorno; conciliación sólo bajo disparo. &
      \textbf{Arrastre}: credenciales confinadas; conciliación continua;
      un agente por entorno. \\
    Mutabilidad &
      \textbf{Mutable}: cambios rápidos y baratos; deriva acumulativa; difícil
      trazar qué versión corre. &
      \textbf{Inmutable}: sin deriva; reemplazo trazable; obliga a externalizar
      todo estado persistente. \\
    Expresión &
      \textbf{Declarativa}: el motor calcula el camino; correcta desde
      cualquier estado de partida. &
      \textbf{Imperativa}: control fino del orden; sólo correcta desde el
      estado de partida previsto. \\
    \bottomrule
  \end{tabularx}
\end{table}

% ==========================================================================
\section{El módulo como unidad de composición}
\label{sec:modulo}
% ==========================================================================

El término \enquote{módulo} tiene en IaC un sentido preciso: es la unidad
reutilizable que encapsula un conjunto de recursos tras una interfaz de
entradas y salidas. Es el mecanismo que permite que una organización no repita
la misma configuración en cada equipo y, sobre todo, que las decisiones
correctas queden incorporadas por omisión.

\begin{figure}[ht]
  \centering
  \begin{tikzpicture}
    \node[draw=capa, very thick, fill=capa!8, rounded corners=3pt,
      minimum width=5.6cm, minimum height=3.2cm] (mod) {};
    \node[anchor=north, font=\footnotesize\bfseries, text=capa!60!black]
      at ([yshift=-3mm]mod.north) {Módulo \texttt{red-privada v2.3.1}};
    \node[bloque, fill=white, font=\scriptsize, text width=44mm,
      anchor=north] at ([yshift=-9mm]mod.north)
      {\textbf{Cuerpo}: recursos declarados,\\ valores por omisión,
       validaciones,\\ etiquetas obligatorias};

    \node[bloque, fill=blue!8, text width=27mm, font=\scriptsize,
      left=16mm of mod] (in)
      {\textbf{Entradas}\\ región, CIDR,\\ número de zonas,\\ etiquetas};
    \node[bloque, fill=green!10, text width=27mm, font=\scriptsize,
      right=16mm of mod] (out)
      {\textbf{Salidas}\\ identificador de red,\\ subredes,\\ tabla de rutas};
    \node[bloque, fill=trans!12, text width=40mm, font=\scriptsize,
      below=9mm of mod] (reg)
      {\textbf{Registro versionado}\\ publica, firma y resuelve versiones};

    \draw[flecha] (in) -- (mod);
    \draw[flecha] (mod) -- (out);
    \draw[flecha] (reg) -- (mod);

    \node[bloque, fill=gray!8, text width=25mm, font=\scriptsize,
      above left=8mm and 8mm of mod.north] (c1) {Composición\\
      \emph{entorno de pruebas}};
    \node[bloque, fill=gray!8, text width=25mm, font=\scriptsize,
      above right=8mm and 8mm of mod.north] (c2) {Composición\\
      \emph{entorno de producción}};
    \draw[flecha, gray] (c1) -- (mod.north west);
    \draw[flecha, gray] (c2) -- (mod.north east);
  \end{tikzpicture}
  \caption{Anatomía de un módulo. Las mismas entradas con valores distintos
  producen entornos comparables; el registro versionado permite que una
  corrección se propague de forma controlada. Elaboración propia.}
  \label{fig:modulo}
\end{figure}

Tres propiedades hacen que un módulo cumpla su función. Debe tener una
\textbf{interfaz estrecha}: cuantas menos entradas expone, más decisiones fija
y más valor aporta. Debe estar \textbf{versionado} con semántica explícita,
para que quien lo consume elija cuándo adoptar un cambio. Y debe delimitar un
\textbf{radio de explosión} razonable: un módulo que agrupa demasiados
recursos obliga a que cada cambio pequeño planifique sobre un conjunto grande,
lo que alarga el ciclo y aumenta el riesgo de cada aplicación.

% ==========================================================================
\section{Matriz de estructura de dependencia (DSM)}
\label{sec:dsm}
% ==========================================================================

Esta sección corresponde al elemento 2 del índice de una hoja de ruta
tecnológica. Las secciones anteriores han enumerado las piezas del sistema,
pero enumerarlas no basta: lo que decide el orden de una implantación, el
riesgo de cada paso y la conveniencia de agrupar unos trabajos y separar otros
son las dependencias entre ellas. La matriz de estructura de dependencia, o
DSM \parencite{steward1981,eppinger2012}, es el instrumento que las hace
visibles y operables, porque permite reordenar filas y columnas hasta que la
propia matriz muestra la secuencia.

Se construyen dos matrices, porque la hoja de ruta necesita responder a dos
preguntas distintas. La primera es de arquitectura: qué componente del sistema
necesita a qué otro, lo que determina qué se puede diseñar y desplegar por
separado. La segunda es de cartera: qué proyecto de la hoja de ruta necesita
que otro esté terminado, lo que determina el calendario y el orden de
inversión.

\subsec{01. Convención de lectura}

Ambas matrices son cuadradas y llevan los mismos elementos en filas y en
columnas, en el mismo orden. Se adopta la convención de entradas por fila:

\begin{quote}
Una marca en la celda de la fila $i$ y la columna $j$ significa que el
elemento $i$ \textbf{necesita algo del} elemento $j$. Se lee, por tanto,
recorriendo la fila $i$ de izquierda a derecha: las marcas de esa fila son
todo lo que $i$ recibe. Recorriendo en cambio la columna $j$ de arriba abajo
se obtiene todo lo que $j$ entrega.
\end{quote}

De esa convención se derivan las tres lecturas que se usan después. La
\textbf{suma de una fila} mide cuántas condiciones previas tiene un elemento,
es decir, su dificultad de arranque. La \textbf{suma de una columna} mide a
cuántos elementos afecta, es decir, el alcance del daño si falla o se retrasa.
Y una marca \textbf{por encima de la diagonal}, una vez ordenada la matriz,
señala una realimentación: un elemento que necesita algo que en el orden
propuesto todavía no existe, lo que obliga a iterar.

La diagonal no se marca: ningún elemento se necesita a sí mismo. Las celdas
diagonales se sombrean únicamente para orientar la vista.

\subsec{02. La DSM de arquitectura}

Los elementos son los trece componentes ya descritos: las siete capas de la
Figura~\ref{fig:pila} y los seis módulos transversales del
Cuadro~\ref{tab:transversales}. El Cuadro~\ref{tab:dsmelem} fija los códigos
que usa la matriz y, para cada componente, el tipo de dependencia que
introduce.

\begin{table}[ht]
  \centering
  \footnotesize
  \caption{Los trece componentes de la DSM de arquitectura y la dependencia
  que cada uno introduce en el sistema.}
  \label{tab:dsmelem}
  \begin{tabularx}{\textwidth}{@{}l>{\raggedright\arraybackslash}p{0.235\textwidth}
    >{\raggedright\arraybackslash}X@{}}
    \toprule
    \textbf{Cód.} & \textbf{Componente} & \textbf{Qué necesita de los demás} \\
    \midrule
    L0 & Sustrato del proveedor & Nada del sistema: es su objeto. \\
    L1 & Aprovisionamiento & La interfaz de L0, la especificación versionada,
         el registro de estado, las credenciales y el veredicto de la
         política. \\
    L2 & Imágenes & Sustrato donde construir, definiciones versionadas y
         credenciales del registro de artefactos. \\
    L3 & Configuración & Máquinas ya creadas por L1 a partir de imágenes de
         L2. \\
    L4 & Orquestación & El agrupamiento aprovisionado por L1 y las imágenes de
         contenedor de L2. \\
    L5 & Entrega & El motor de L1, el destino de L4, el repositorio, las
         políticas y las pruebas que ejecuta en cada paso. \\
    L6 & Plataforma & El catálogo que ejecuta sobre L1, la canalización de L5,
         las políticas que codifica y las métricas que muestra. \\
    M1 & Control de versiones & Nada: es la raíz del sistema. \\
    M2 & Gestión de estado & Almacenamiento remoto en L0 y claves de cifrado
         del módulo de secretos. \\
    M3 & Gestión de secretos & El servicio de custodia y cifrado de L0. \\
    M4 & Política como código & El plan que produce L1 y el repositorio donde
         viven las reglas. \\
    M5 & Pruebas & Entornos efímeros que crea L1 y el repositorio del que
         parten. \\
    M6 & Observabilidad & Telemetría de L0, ejecuciones de sólo lectura de L1,
         registros de L5 y el estado custodiado por M2. \\
    \bottomrule
  \end{tabularx}
\end{table}

La Figura~\ref{fig:dsmarq} presenta la matriz en el orden natural de
exposición, capas primero y módulos después. Es el orden en que se ha descrito
el sistema, no el orden en que conviene construirlo, y por eso las marcas
aparecen repartidas a ambos lados de la diagonal.

\begin{figure}[ht]
  \centering
  \begin{tikzpicture}[x=6.4mm, y=6.4mm, font=\scriptsize]
    \def\n{13}
    % ---- marcas: fila/columna (la fila necesita algo de la columna)
    \def\marcas{2/1,2/8,2/9,2/10,2/11,
                3/1,3/8,3/10,
                4/2,4/3,4/8,4/10,
                5/2,5/3,5/8,5/10,
                6/2,6/5,6/8,6/10,6/11,6/12,
                7/2,7/6,7/8,7/11,7/13,
                9/1,9/10,
                10/1,
                11/2,11/8,
                12/2,12/8,
                13/1,13/2,13/6,13/9}
    % ---- fondo de la rejilla y diagonal
    \fill[gray!4] (0,0) rectangle (\n,-\n);
    \foreach \i in {1,...,13}
      \fill[gray!35] (\i-1,-\i+1) rectangle (\i,-\i);
    % ---- marcas
    \foreach \r/\c in \marcas
      \fill[dsmarq] (\c-1+0.14,-\r+1-0.14) rectangle (\c-0.14,-\r+0.14);
    % ---- rejilla
    \draw[gray!55, thin, step=1] (0,0) grid (\n,-\n);
    \draw[thick] (0,0) rectangle (\n,-\n);
    % ---- separación entre capas y módulos
    \draw[thick, gray!70] (7,0) -- (7,-\n);
    \draw[thick, gray!70] (0,-7) -- (\n,-7);
    % ---- etiquetas
    \foreach \i/\lab in {1/L0,2/L1,3/L2,4/L3,5/L4,6/L5,7/L6,
                         8/M1,9/M2,10/M3,11/M4,12/M5,13/M6} {
      \node[anchor=east] at (-0.2,-\i+0.5) {\lab};
      \node[anchor=west, rotate=90] at (\i-0.5,0.2) {\lab};
    }
    % ---- sumas
    \foreach \i/\v in {1/0,2/5,3/3,4/4,5/4,6/6,7/5,8/0,9/2,10/1,11/2,12/2,13/4}
      \node[anchor=west, font=\scriptsize\bfseries] at (\n+0.35,-\i+0.5) {\v};
    \foreach \i/\v in {1/5,2/7,3/2,4/0,5/1,6/2,7/0,8/8,9/2,10/6,11/3,12/1,13/1}
      \node[anchor=north, font=\scriptsize\bfseries] at (\i-0.5,-\n-0.25) {\v};
    \node[anchor=west, font=\scriptsize\itshape, align=left]
      at (\n+0.35,0.7) {entradas};
    \node[anchor=north, font=\scriptsize\itshape] at (\n+1.1,-\n-0.25)
      {salidas};
  \end{tikzpicture}
  \caption{DSM de arquitectura en el orden natural de exposición. Una marca en
  la fila $i$ y la columna $j$ indica que el componente $i$ necesita algo del
  componente $j$. La columna de la derecha suma las entradas de cada fila y la
  banda inferior suma las salidas de cada columna. Elaboración propia.}
  \label{fig:dsmarq}
\end{figure}

\subsec{03. Qué dice la matriz sin reordenarla}

Antes de tocar el orden, las sumas ya dicen tres cosas.

\textbf{Dos de los tres concentradores de riesgo son módulos transversales, no
capas.} Las columnas con más salidas son M1, el control de versiones, con
ocho; L1, el aprovisionamiento, con siete; y M3, la gestión de secretos, con
seis. Es un resultado con consecuencias prácticas: una interrupción del
repositorio detiene ocho componentes y una del custodio de secretos detiene
seis, mientras que la caída de la capa de configuración, L3, con cero salidas,
no detiene a ninguno. La disponibilidad que se exige a cada pieza debería
seguir ese orden, y con frecuencia no lo sigue, porque la atención se
concentra en el motor, que es sólo el segundo de los tres.

\textbf{Las capas superiores son las caras de construir, no las de operar.} Las
filas con más entradas son L5, la entrega, con seis, y L1 y L6 con cinco. Un
proyecto que empiece por la interfaz de autoservicio está aceptando cinco
condiciones previas simultáneas; empezar por ahí es la causa habitual de que
un portal interno quede a medias.

\textbf{Hay dos elementos sin ninguna entrada}, L0 y M1, que son por tanto los
únicos puntos de arranque posibles de cualquier secuencia.

\subsec{04. Particionado: la secuencia de implantación}

Reordenar filas y columnas con la misma permutación no cambia la información
de la matriz, sólo la presentación. El particionado busca la permutación que
lleva el mayor número posible de marcas por debajo de la diagonal, porque una
matriz triangular inferior describe un orden en el que cada elemento se
construye cuando todo lo que necesita ya existe. La
Figura~\ref{fig:dsmarqord} recoge el resultado.

\begin{figure}[ht]
  \centering
  \begin{tikzpicture}[x=6.4mm, y=6.4mm, font=\scriptsize]
    \def\n{13}
    % Orden: M1 L0 M3 M2 L2 L1 M4 M5 L3 L4 L5 M6 L6
    % Posición nueva: 1=M1 2=L0 3=M3 4=M2 5=L2 6=L1 7=M4 8=M5 9=L3 10=L4
    %                 11=L5 12=M6 13=L6
    \def\marcas{3/2, 4/2,4/3, 5/1,5/2,5/3, 6/1,6/2,6/3,6/4, 7/1,7/6,
                8/1,8/6, 9/1,9/3,9/5,9/6, 10/1,10/3,10/5,10/6,
                11/1,11/3,11/6,11/7,11/8,11/10,
                12/2,12/4,12/6,12/11, 13/1,13/6,13/7,13/11,13/12}
    \fill[gray!4] (0,0) rectangle (\n,-\n);
    \foreach \i in {1,...,13}
      \fill[gray!35] (\i-1,-\i+1) rectangle (\i,-\i);
    % ---- marca de realimentación (por encima de la diagonal)
    \fill[acento] (6.14,-5.14) rectangle (6.86,-5.86);
    \foreach \r/\c in \marcas
      \fill[dsmarq] (\c-1+0.14,-\r+1-0.14) rectangle (\c-0.14,-\r+0.14);
    \draw[gray!55, thin, step=1] (0,0) grid (\n,-\n);
    \draw[thick] (0,0) rectangle (\n,-\n);
    % ---- escalonado de las olas de implantación, sobre la rejilla
    \foreach \a/\b in {0/2, 2/3, 3/5, 5/7, 7/10, 10/11, 11/12, 12/13}
      \draw[gray!75!black, line width=1pt] (\a,-\a) -- (\b,-\a) -- (\b,-\b);
    % ---- bloque acoplado L1-M4
    \draw[acento, very thick] (5,-5) rectangle (7,-7);
    \foreach \i/\lab in {1/M1,2/L0,3/M3,4/M2,5/L2,6/L1,7/M4,
                         8/M5,9/L3,10/L4,11/L5,12/M6,13/L6} {
      \node[anchor=east] at (-0.2,-\i+0.5) {\lab};
      \node[anchor=west, rotate=90] at (\i-0.5,0.2) {\lab};
    }
    \node[anchor=west, font=\scriptsize\itshape, text=acento]
      at (7.4,-6) {bloque acoplado};
    \draw[acento, thick] (7.05,-5.5) -- (7.35,-6);
  \end{tikzpicture}
  \caption{DSM de arquitectura particionada. En el orden M1, L0, M3, M2, L2,
  L1, M4, M5, L3, L4, L5, M6, L6 todas las dependencias caen por debajo de la
  diagonal salvo una, la que L1 tiene de M4, marcada en rojo. El escalonado
  gris separa las olas: los elementos de una misma ola no dependen entre sí y
  pueden abordarse en paralelo. Elaboración propia.}
  \label{fig:dsmarqord}
\end{figure}

El resultado es una matriz triangular inferior con una sola excepción, y esa
excepción es informativa. El escalonado define ocho olas, y los elementos que
comparten ola pueden abordarse en paralelo:

\begin{enumerate}
  \item \textbf{M1 y L0}: el repositorio y la cuenta del proveedor. Sin
    entradas, son el arranque obligado.
  \item \textbf{M3}: la gestión de secretos. Aparece la segunda, antes que
    cualquier capa, porque seis componentes la necesitan y ninguno puede
    autenticarse sin ella. Ponerla al final, que es lo habitual cuando se
    trata como un endurecimiento posterior, obliga a rehacer todo lo anterior.
  \item \textbf{M2 y L2}: la custodia del registro de estado y la construcción
    de imágenes, que no dependen la una de la otra.
  \item \textbf{L1 y M4}: el bloque acoplado, del que se habla enseguida.
  \item \textbf{M5, L3 y L4}: las pruebas, la configuración y la orquestación,
    los tres apoyados en el motor ya operativo.
  \item \textbf{L5}: la entrega continua, que necesita el motor, la
    orquestación, las políticas y las pruebas a la vez.
  \item \textbf{M6}: la observabilidad, que sólo puede medir lo que ya
    funciona.
  \item \textbf{L6}: la interfaz de autoservicio, la última porque envuelve a
    todo lo demás.
\end{enumerate}

\textbf{El bloque acoplado y cómo se rompe.} L1 necesita de M4 el veredicto de
la política antes de aplicar un cambio, y M4 necesita de L1 el plan sobre el
que pronunciarse. En una DSM eso es un ciclo, y un ciclo obliga a iterar o a
descomponer. La descomposición correcta consiste en partir L1 en dos
subprocesos, uno que produce el plan y otro que lo aplica, con la evaluación
de política entre ambos. Es exactamente lo que hace el diagrama SD1 del modelo
OPM de la sección~\ref{sec:opm}: al separar \emph{Planificación},
\emph{Validación} y \emph{Aplicación} en subprocesos distintos, el ciclo
desaparece y la secuencia queda lineal. La DSM detecta el acoplamiento y el
modelo OPM propone la partición que lo elimina, lo que ilustra por qué los
elementos 2 y 3 de la hoja de ruta se sostienen mutuamente.

\subsec{05. La DSM de la cartera de proyectos}

La segunda matriz cambia de objeto. Sus elementos ya no son componentes del
sistema sino los ocho proyectos con que la empresa pasa de la plataforma
actual a la plataforma de 2029. Tres de ellos se corresponden con los vectores
tecnológicos del Cuadro~\ref{tab:vectores}, que la sección~\ref{sec:comp}
utiliza para construir el gráfico vectorial; los otros cinco son los
habilitadores sin los cuales aquéllos no se sostienen. El
Cuadro~\ref{tab:dsmproy} los enumera.

\begin{table}[ht]
  \centering
  \footnotesize
  \caption{Los ocho proyectos de la cartera, su correspondencia con los
  vectores tecnológicos del Cuadro~\ref{tab:vectores} y los componentes de la
  DSM de arquitectura sobre los que actúan. Las filas con la columna
  \emph{vector} en blanco son proyectos habilitadores, que no desplazan por sí
  solos ninguna coordenada del gráfico vectorial.}
  \label{tab:dsmproy}
  \begin{tabularx}{\textwidth}{@{}l>{\raggedright\arraybackslash}p{0.245\textwidth}
    >{\raggedright\arraybackslash}X l l@{}}
    \toprule
    \textbf{Cód.} & \textbf{Proyecto} & \textbf{Alcance} & \textbf{Vector} &
      \textbf{Comp.} \\
    \midrule
    R1 & Cierre de cobertura de codificación & Inventariar los recursos
         creados a mano e importarlos al motor hasta llevar la cobertura del
         55 al 86 por ciento. & & L1 \\
    R2 & Custodia del registro de estado & Estado remoto, cifrado, versionado
         y con bloqueo de concurrencia. & & M2 \\
    R3 & Catálogo de módulos certificados & Módulos versionados, firmados y
         publicados en un registro interno, con interfaz estrecha. & P1 &
         L1, L6 \\
    R4 & Política como código & Reglas de la organización evaluadas sobre el
         plan antes de aplicarlo. & P1 & M4 \\
    R5 & Malla de pruebas y entornos efímeros & Análisis estático, validación
         del plan y despliegue efímero verificado. & & M5 \\
    R6 & Sustitución del motor & Migración a la bifurcación abierta con
         estado gestionado propio. & P2 & L1 \\
    R7 & Interfaz de autoservicio & Portal que expone el catálogo al equipo
         de producto. & P3 & L6 \\
    R8 & Observabilidad de la plataforma & Deriva, coste, tiempo de ciclo y
         cumplimiento como series temporales. & & M6 \\
    \bottomrule
  \end{tabularx}
\end{table}

Las dependencias entre proyectos no son las mismas que entre componentes,
porque un proyecto puede necesitar de otro un resultado parcial y no el
componente entero. La Figura~\ref{fig:dsmproy} distingue por eso dos
intensidades: dependencia \emph{necesaria}, cuando el proyecto no puede
empezar sin el resultado del otro, y dependencia \emph{conveniente}, cuando
puede empezar pero a un coste o un riesgo mayores.

\begin{figure}[ht]
  \centering
  \begin{tikzpicture}[x=8.4mm, y=8.4mm, font=\scriptsize]
    \def\n{8}
    \def\fuerte{1/2, 3/1, 4/3, 5/3, 6/2, 6/3, 6/5, 7/3, 7/4, 8/1}
    \def\debil{3/2, 4/1, 5/2, 7/6, 8/2}
    \fill[gray!4] (0,0) rectangle (\n,-\n);
    \foreach \i in {1,...,8}
      \fill[gray!35] (\i-1,-\i+1) rectangle (\i,-\i);
    \foreach \r/\c in \fuerte
      \fill[dsmpro] (\c-1+0.16,-\r+1-0.16) rectangle (\c-0.16,-\r+0.16);
    \foreach \r/\c in \debil
      \draw[dsmpro, thick, fill=dsmpro!20]
        (\c-1+0.16,-\r+1-0.16) rectangle (\c-0.16,-\r+0.16);
    \draw[gray!55, thin, step=1] (0,0) grid (\n,-\n);
    \draw[thick] (0,0) rectangle (\n,-\n);
    \foreach \i/\lab in {1/R1,2/R2,3/R3,4/R4,5/R5,6/R6,7/R7,8/R8} {
      \node[anchor=east] at (-0.2,-\i+0.5) {\lab};
      \node[anchor=south] at (\i-0.5,0.1) {\lab};
    }
  \end{tikzpicture}
  \caption{DSM de la cartera de proyectos en el orden de codificación R1 a R8.
  Marca llena: el proyecto de la fila no puede empezar sin el resultado del
  proyecto de la columna. Marca hueca: puede empezar, pero a mayor coste o
  riesgo. Elaboración propia.}
  \label{fig:dsmproy}
\end{figure}

\subsec{06. Secuenciación de la cartera}

Aplicado el mismo particionado, la matriz de la cartera resulta ser
completamente triangular inferior: no hay ningún ciclo, de modo que existe una
secuencia que no obliga a rehacer nada. La Figura~\ref{fig:dsmproyord} la
recoge y la traduce en seis olas.

\begin{figure}[ht]
  \centering
  \begin{tikzpicture}[x=8.4mm, y=8.4mm, font=\scriptsize]
    \def\n{8}
    % Orden: R2 R1 R3 R8 R4 R5 R6 R7
    % Posición: 1=R2 2=R1 3=R3 4=R8 5=R4 6=R5 7=R6 8=R7
    \def\fuerte{2/1, 3/2, 4/2, 5/3, 6/3, 7/1, 7/3, 7/6, 8/3, 8/5}
    \def\debil{3/1, 4/1, 5/2, 6/1, 8/7}
    \fill[gray!4] (0,0) rectangle (\n,-\n);
    \foreach \i in {1,...,8}
      \fill[gray!35] (\i-1,-\i+1) rectangle (\i,-\i);
    \foreach \r/\c in \fuerte
      \fill[dsmpro] (\c-1+0.16,-\r+1-0.16) rectangle (\c-0.16,-\r+0.16);
    \foreach \r/\c in \debil
      \draw[dsmpro, thick, fill=dsmpro!20]
        (\c-1+0.16,-\r+1-0.16) rectangle (\c-0.16,-\r+0.16);
    \draw[gray!55, thin, step=1] (0,0) grid (\n,-\n);
    \draw[thick] (0,0) rectangle (\n,-\n);
    % ---- escalonado de las olas, sobre la rejilla
    \foreach \a/\b in {0/1, 1/2, 2/4, 4/6, 6/7, 7/8}
      \draw[gray!75!black, line width=1pt] (\a,-\a) -- (\b,-\a) -- (\b,-\b);
    \foreach \i/\lab in {1/R2,2/R1,3/R3,4/R8,5/R4,6/R5,7/R6,8/R7} {
      \node[anchor=east] at (-0.2,-\i+0.5) {\lab};
      \node[anchor=south] at (\i-0.5,0.1) {\lab};
    }
    \foreach \y/\ola in {0.5/{ola 1}, 1.5/{ola 2}, 3/{ola 3}, 5/{ola 4},
                         6.5/{ola 5}, 7.5/{ola 6}}
      \node[anchor=west, font=\scriptsize\itshape, text=gray!55!black]
        at (\n+0.35,-\y) {\ola};
  \end{tikzpicture}
  \caption{DSM de la cartera particionada. Todas las marcas quedan por debajo
  de la diagonal: la cartera admite una secuencia sin iteración. R3 y R8 caen
  en la misma ola porque no dependen entre sí, y lo mismo ocurre con R4 y R5.
  Elaboración propia.}
  \label{fig:dsmproyord}
\end{figure}

La secuencia resultante es R2; R1; R3 y R8 en paralelo; R4 y R5 en paralelo;
R6; y por último R7. Tres consecuencias merecen comentario.

\textbf{La custodia del estado va primero, y sola.} R2 no depende de nada y
casi todo depende de él. Es el único proyecto de la cartera que no admite ser
pospuesto, porque importar recursos a un registro de estado sin bloqueo ni
cifrado, que es lo que haría R1 sin él, multiplica el trabajo por dos: primero
se importa y después se migra.

\textbf{La observabilidad no es el remate, es la ola 3.} R8 aparece pronto
porque sin él las figuras de mérito del Cuadro~\ref{tab:fom} no son medibles,
y sin medida las posiciones del gráfico vectorial de la
sección~\ref{sec:comp} siguen siendo estimaciones. Situarlo al final, que es
la tentación natural, deja la hoja de ruta sin instrumento de verificación
justo durante los años en que se ejecuta.

\textbf{La matriz confirma el orden discutido del catálogo y el motor.} El
apartado 05 de la sección~\ref{sec:comp} defiende en prosa que el catálogo de
módulos (R3) debe construirse antes que la sustitución del motor (R6), y
reconoce que es la decisión más discutible de aquella sección. La DSM lo
convierte en una constatación estructural: R6 tiene marca necesaria en las
columnas R2, R3 y R5, de modo que en cualquier orden admisible R6 queda por
detrás de los tres. Invertir el orden no es una preferencia distinta, es una
marca por encima de la diagonal, esto es, una iteración forzada.

\subsec{07. Riesgos que la matriz hace visibles}

Cruzar las dos matrices identifica dónde conviene poner las reservas de
esfuerzo, que es el uso que la gestión del portafolio hace de la DSM.

\begin{table}[!ht]
  \centering
  \footnotesize
  \caption{Riesgos derivados de la estructura de dependencias y respuesta
  prevista.}
  \label{tab:dsmriesgo}
  \begin{tabularx}{\textwidth}{@{}>{\raggedright\arraybackslash}p{0.235\textwidth}
    >{\raggedright\arraybackslash}X
    >{\raggedright\arraybackslash}p{0.30\textwidth}@{}}
    \toprule
    \textbf{Origen en la matriz} & \textbf{Riesgo} & \textbf{Respuesta} \\
    \midrule
    M1 con ocho salidas & Una indisponibilidad del repositorio detiene ocho
      componentes; su compromiso los compromete a todos. &
      Réplica del repositorio y firma obligatoria de confirmaciones. \\
    M3 con seis salidas y ola temprana & Si la gestión de secretos llega tarde,
      lo construido antes queda con credenciales incrustadas. &
      Adelantarla a la ola 2 de la arquitectura, antes que cualquier capa
      intermedia. \\
    Bloque acoplado L1--M4 & La política evaluada después de aplicar convierte
      el incumplimiento en un hecho consumado. &
      Partir el motor en planificación y aplicación, con la evaluación
      intercalada, según el diagrama SD1 de la sección~\ref{sec:opm}. \\
    R2 con cinco entradas en su columna & Un retraso en la custodia del estado
      desplaza toda la cartera. &
      Ola 1 en solitario, con criterio de terminación explícito antes de
      abrir R1. \\
    R3 con cuatro entradas en su columna & El catálogo es el cuello de botella
      de la cartera: R4, R5, R6 y R7 lo esperan. &
      Publicar el catálogo por incrementos y abrir R4 y R5 contra los primeros
      módulos, sin esperar al catálogo completo. \\
    L3 y L6 con cero salidas & Ningún componente los espera, luego son los
      candidatos naturales a recortar si el calendario se comprime. &
      Registrar el recorte como decisión consciente, no como omisión. \\
    \bottomrule
  \end{tabularx}
\end{table}

\subsec{08. Límites de las dos matrices}

Tres reservas. La primera es que ambas matrices son \textbf{binarias en el
tipo}: registran que hay dependencia, no de qué naturaleza es. Una DSM más
elaborada distinguiría dependencias de información, de material, de energía y
espaciales, y la de este sistema sería casi enteramente de información, lo que
resta valor a la distinción.

La segunda es que la DSM de arquitectura describe una \textbf{implantación
completa desde cero}. Una organización que ya opera parte del sistema, que es
el caso de la empresa de referencia, no recorre las olas en orden sino que
entra en la secuencia por el punto en que se encuentra; el orden sigue siendo
válido para los componentes que faltan.

La tercera es que la matriz de cartera \textbf{no incorpora el calendario ni
el coste}. Dice qué debe ir antes que qué, no cuánto dura cada cosa ni cuántos
equipos caben en paralelo. Convertir estas olas en fechas exige el modelo
financiero, el elemento 8 del índice, que este documento declara fuera de
alcance.

% ==========================================================================
\section{Modelo de objetos y procesos (OPM)}
\label{sec:opm}
% ==========================================================================

\subsec{01. Por qué OPM y cómo se ha construido el modelo}

Esta sección corresponde al elemento 3 del índice de una hoja de ruta
tecnológica. La metodología de objetos y procesos, normalizada como ISO/PAS
19450:2015 y después como ISO 19450:2024
\parencite{iso19450:2015,iso19450:2024,dori2002,dori2016}, presenta dos rasgos
que la hacen adecuada aquí. El primero es que usa un único tipo de diagrama
para representar a la vez estructura y comportamiento, lo que evita el
problema habitual de mantener sincronizados un diagrama de bloques y un
diagrama de actividad. El segundo es su carácter bimodal: cada diagrama tiene
una traducción a lenguaje natural controlado, el OPL, de modo que el modelo se
verifica leyéndolo, y un enlace mal puesto se manifiesta como una frase falsa.

Las herramientas autorizadas para este elemento son OPCloud
\parencite{opcloud} y OPCAT \parencite{opcat}, y el modelo está enunciado en
el vocabulario cerrado que ambas admiten: cada entidad, cada enlace y cada
frase del OPL de los apartados siguientes se corresponden con un elemento de
la norma y con ningún otro. Las figuras de esta sección son la composición
tipográfica de ese mismo modelo, dibujada con la notación de la norma para que
quepa en la caja del documento y pueda referirse desde el texto.

\subsec{02. Vocabulario y notación empleados}

OPM tiene un vocabulario cerrado, y la mayor parte de los errores de modelado
consisten en usar un símbolo por otro. El Cuadro~\ref{tab:opmnotacion} fija el
que se usa en las tres figuras siguientes.

\begin{table}[ht]
  \centering
  \footnotesize
  \caption{Vocabulario de OPM empleado en las Figuras \ref{fig:opmsd},
  \ref{fig:opmsd2} y \ref{fig:opmsd1}, con la frase OPL que genera cada
  símbolo.}
  \label{tab:opmnotacion}
  \begin{tabularx}{\textwidth}{@{}>{\raggedright\arraybackslash}p{0.215\textwidth}
    >{\raggedright\arraybackslash}p{0.305\textwidth}
    >{\raggedright\arraybackslash}X@{}}
    \toprule
    \textbf{Elemento} & \textbf{Símbolo} & \textbf{Frase OPL} \\
    \midrule
    \multicolumn{3}{@{}l}{\textit{Entidades}} \\
    Objeto & Rectángulo de contorno verde & \\
    Proceso & Elipse de contorno azul & \\
    Estado & Rectángulo redondeado dentro del objeto &
      \emph{O} puede estar \emph{s1}, \emph{s2} o \emph{s3}. \\
    Cosa física & Con sombra proyectada & \emph{C} es física. \\
    Cosa ambiental & Contorno discontinuo & \emph{C} es ambiental. \\
    \midrule
    \multicolumn{3}{@{}l}{\textit{Enlaces procedimentales}} \\
    Agente & Línea acabada en círculo relleno &
      \emph{O} maneja \emph{P}. \\
    Instrumento & Línea acabada en círculo hueco &
      \emph{P} requiere \emph{O}. \\
    Efecto & Flecha de doble punta &
      \emph{P} afecta a \emph{O}. \\
    Consumo & Flecha del objeto al proceso &
      \emph{P} consume \emph{O}. \\
    Resultado & Flecha del proceso al objeto &
      \emph{P} produce \emph{O}. \\
    Entrada y salida & Flechas entre estados concretos y el proceso &
      \emph{P} cambia \emph{O} de \emph{s1} a \emph{s2}. \\
    \midrule
    \multicolumn{3}{@{}l}{\textit{Enlaces estructurales}} \\
    Agregación & Triángulo relleno &
      \emph{A} consta de \emph{B}. \\
    Generalización & Triángulo hueco &
      \emph{B} es un \emph{A}. \\
    Exhibición & Triángulo dentro de triángulo &
      \emph{A} exhibe \emph{B}. \\
    \midrule
    \multicolumn{3}{@{}l}{\textit{Mecanismos de refinamiento}} \\
    Ampliación en detalle & El proceso ampliado contiene a sus subprocesos,
      ordenados de arriba abajo por precedencia temporal &
      \emph{P} se amplía en \emph{P1}, \emph{P2} y \emph{P3}. \\
    Despliegue & Diagrama aparte con la jerarquía estructural de un objeto &
      (las de sus enlaces estructurales) \\
    \bottomrule
  \end{tabularx}
\end{table}

Dos precisiones sobre el uso de estos símbolos, porque son las dos que con más
frecuencia se incumplen. La primera: la precedencia temporal entre subprocesos
de una ampliación se expresa \textbf{únicamente} por la posición vertical, y
no dibujando flechas entre ellos; una flecha entre dos procesos sería un
enlace de invocación, que significa otra cosa. La segunda: cuando el diagrama
muestra los estados de un objeto y el proceso transforma uno concreto en otro,
los enlaces deben partir y llegar \textbf{al estado} y no al objeto entero; el
enlace de efecto al objeto entero sólo es correcto cuando el cambio de estado
no se especifica.

\subsec{03. Diagrama de sistema (SD)}

El diagrama de sistema recoge el proceso principal, quién lo maneja, qué
instrumentos necesita y qué objetos transforma o produce. Los estados de la
\emph{Infraestructura Desplegada} se declaran aquí, pero las transiciones
concretas entre ellos no se especifican todavía: a este nivel el proceso la
afecta, sin más, y las transiciones aparecen en la ampliación en detalle. El
Cuadro~\ref{tab:entidades} enumera las entidades.

\begin{table}[ht]
  \centering
  \scriptsize
  \caption{Entidades del modelo OPM y diagrama en el que aparecen.}
  \label{tab:entidades}
  \begin{tabularx}{\textwidth}{@{}>{\raggedright\arraybackslash}p{0.245\textwidth}
    >{\raggedright\arraybackslash}p{0.155\textwidth}
    >{\raggedright\arraybackslash}X@{}}
    \toprule
    \textbf{Entidad} & \textbf{Tipo} & \textbf{Papel en el modelo} \\
    \midrule
    Aprovisionamiento de Infraestructura & Proceso & Proceso principal; se
      amplía en detalle en SD1. \\
    Equipo de Plataforma & Objeto físico y ambiental & Agente humano que
      maneja el proceso. \\
    Especificación de Infraestructura & Objeto informático & El código;
      instrumento del proceso. Se despliega en SD2. Estados: borrador,
      validada. \\
    Repositorio Versionado & Objeto informático & Todo del que la
      especificación es parte (SD2). \\
    Módulo Reutilizable, Configuración de Entorno & Objetos informáticos &
      Especializaciones de la especificación (SD2). \\
    Versión & Objeto informático & Atributo exhibido por la especificación
      (SD2). \\
    Motor de IaC & Objeto informático & Instrumento que ejecuta la
      comparación y la conciliación. \\
    Registro de Estado & Objeto informático & Objeto afectado: el proceso lo
      lee y lo reescribe en cada aplicación. \\
    Proveedor de Nube & Objeto físico y ambiental & Instrumento: expone la
      interfaz de programación. \\
    Credenciales & Objeto informático & Instrumento con el que el motor se
      autentica. \\
    Conjunto de Políticas & Objeto informático & Instrumento de la
      validación. \\
    Infraestructura Desplegada & Objeto físico & Operando del sistema.
      Estados: inexistente, conforme, desviada, retirada. \\
    Plan de Cambios & Objeto informático & Resultado de la planificación,
      instrumento de la aplicación. Estados: pendiente, aprobado. \\
    Registro de Auditoría & Objeto informático & Resultado del proceso. \\
    \bottomrule
  \end{tabularx}
\end{table}

\begin{figure}[H]
  \centering
  \begin{tikzpicture}[scale=0.86, transform shape,
                      every node/.style={transform shape}]
    % ---- proceso principal
    \node[opmproc, minimum width=44mm, minimum height=17mm] (P) at (0,0)
      {Aprovisionamiento\\de Infraestructura};

    % ---- agente humano: objeto físico y ambiental
    \node[opmobj, opmamb, opmfis, text width=24mm] (eq) at (0,3.05)
      {Equipo de\\Plataforma};
    \draw[agente] (eq) -- (P);

    % ---- instrumentos
    \node[opmobj] (spec)  at (-5.8, 1.85) {Especificación de\\Infraestructura};
    \node[opmobj] (motor) at (-5.8, 0.35) {Motor de IaC};
    \node[opmobj] (pol)   at (-5.8,-1.15) {Conjunto de\\Políticas};
    \node[opmobj, opmamb, opmfis] (prov) at (5.8, 1.85)
      {Proveedor\\de Nube};
    \node[opmobj] (cred)  at ( 5.8, 0.35) {Credenciales};
    \foreach \nd in {spec,motor,pol,prov,cred} \draw[instr] (\nd) -- (P);

    % ---- objeto afectado
    \node[opmobj] (edo) at (-5.8,-2.75) {Registro\\de Estado};
    \draw[efecto] (edo) -- (P);

    % ---- resultados
    \node[opmobj] (plan) at (5.8,-1.15) {Plan de Cambios};
    \node[opmobj] (aud)  at (5.8,-2.75) {Registro de\\Auditoría};
    \draw[flecha] (P) -- (plan);
    \draw[flecha] (P) -- (aud);

    % ---- operando: objeto físico con sus cuatro estados
    \node[draw=opmobjc, very thick, fill=opmobjc!7, opmfis,
      minimum width=84mm, minimum height=18mm] (infra) at (0,-4.35) {};
    \node[anchor=north, font=\footnotesize\bfseries]
      at ([yshift=-1.6mm]infra.north) {Infraestructura Desplegada};
    \foreach \dx/\st in {-30/inexistente, -10/conforme, 10/desviada,
                          30/retirada}
      \node[opmstate] at ([xshift=\dx mm, yshift=-4.6mm]infra.center) {\st};
    \draw[efecto] (P) -- (infra);
  \end{tikzpicture}
  \caption{Diagrama de sistema (SD). El proceso principal, su agente, sus
  cinco instrumentos, el objeto que afecta, los dos que produce y el operando
  con sus cuatro estados. El \emph{Equipo de Plataforma} y el \emph{Proveedor
  de Nube} llevan contorno discontinuo por ser ambientales y sombra por ser
  físicos; la \emph{Infraestructura Desplegada} lleva sombra por ser física.
  Elaboración propia según \textcite{dori2016} y \textcite{iso19450:2024}.}
  \label{fig:opmsd}
\end{figure}

\subsec{04. El OPL del diagrama de sistema}

Las palabras reservadas del lenguaje van en versalitas los nombres de las
entidades y en redonda las palabras estructurales. Cada frase se corresponde
con un símbolo de la Figura~\ref{fig:opmsd} y con ningún otro.

\begin{quote}\small
\textsc{Equipo de Plataforma} es físico y ambiental.\\
\textsc{Proveedor de Nube} es físico y ambiental.\\
\textsc{Infraestructura Desplegada} es física.\\
\textsc{Infraestructura Desplegada} puede estar inexistente, conforme,
desviada o retirada.\\
\textsc{Equipo de Plataforma} maneja \textsc{Aprovisionamiento de
Infraestructura}.\\
\textsc{Aprovisionamiento de Infraestructura} requiere \textsc{Especificación
de Infraestructura}, \textsc{Motor de IaC}, \textsc{Conjunto de Políticas},
\textsc{Credenciales} y \textsc{Proveedor de Nube}.\\
\textsc{Aprovisionamiento de Infraestructura} afecta a \textsc{Registro de
Estado}.\\
\textsc{Aprovisionamiento de Infraestructura} afecta a \textsc{Infraestructura
Desplegada}.\\
\textsc{Aprovisionamiento de Infraestructura} produce \textsc{Plan de Cambios}
y \textsc{Registro de Auditoría}.
\end{quote}

Leído en voz alta, el conjunto es una descripción funcional completa del nivel
superior del sistema.

\subsec{05. Despliegue de la especificación (SD2)}

La \emph{Especificación de Infraestructura} tiene una estructura interna que
el diagrama de sistema no muestra, porque mezclarla allí saturaría el nivel
superior. OPM prevé para eso el despliegue: un diagrama aparte que expone la
jerarquía estructural de un objeto sin tocar su comportamiento. La
Figura~\ref{fig:opmsd2} lo recoge con las tres relaciones estructurales de la
norma, que son las tres que este objeto necesita.

\begin{figure}[H]
  \centering
  \begin{tikzpicture}[scale=0.92, transform shape,
                      every node/.style={transform shape}]
    \node[opmobj, text width=32mm] (spec) at (0,0)
      {Especificación de\\Infraestructura};

    % ---- agregación: el repositorio consta de la especificación
    \node[opmobj, text width=30mm] (repo) at (0,2.3)
      {Repositorio\\Versionado};
    \draw[very thick] (repo.south) -- (spec.north);
    \coordinate (ca) at ($(repo.south)!0.45!(spec.north)$);
    \opmagg{0}{ca}
    \node[etiq, anchor=west] at ([xshift=3mm]ca)
      {agregación: \emph{consta de}};

    % ---- exhibición: la especificación exhibe el atributo Versión
    \node[opmobj, text width=17mm] (ver) at (5.9,0) {Versión};
    \draw[very thick] (spec.east) -- (ver.west);
    \coordinate (ce) at ($(spec.east)!0.40!(ver.west)$);
    \opmexh{90}{ce}
    \node[etiq, anchor=south] at ([yshift=2.6mm]$(spec.east)!0.62!(ver.west)$)
      {exhibición: \emph{exhibe}};

    % ---- generalización: dos especializaciones
    \node[opmobj, text width=28mm] (mod) at (-3.0,-2.8)
      {Módulo\\Reutilizable};
    \node[opmobj, text width=28mm] (cfg) at ( 3.0,-2.8)
      {Configuración\\de Entorno};
    \coordinate (cg) at (0,-1.15);
    \draw[very thick] (spec.south) -- (0,-1.75);
    \draw[very thick] (mod.north) |- (0,-1.75) -| (cfg.north);
    \opmgen{0}{cg}
    \node[etiq, anchor=west] at ([xshift=3mm]cg)
      {generalización: \emph{es un}};
  \end{tikzpicture}
  \caption{SD2: despliegue del objeto \emph{Especificación de
  Infraestructura}. Agregación con el repositorio que la contiene,
  generalización en sus dos especializaciones y exhibición del atributo que la
  versiona. La base de cada triángulo se apoya en el todo, el general o el
  exhibidor, y el vértice apunta a la parte, la especialización o el atributo.
  Elaboración propia.}
  \label{fig:opmsd2}
\end{figure}

\begin{quote}\small
\textsc{Repositorio Versionado} consta de \textsc{Especificación de
Infraestructura}.\\
\textsc{Especificación de Infraestructura} exhibe \textsc{Versión}.\\
\textsc{Módulo Reutilizable} es una \textsc{Especificación de
Infraestructura}.\\
\textsc{Configuración de Entorno} es una \textsc{Especificación de
Infraestructura}.
\end{quote}

Conviene subrayar la diferencia entre la tercera frase y una que se le parece
mucho pero significa otra cosa. \emph{Es un} declara una especialización, esto
es, una clase de objeto distinta que hereda las propiedades de la general.
\emph{Puede estar} declara un estado, esto es, una situación transitoria del
mismo objeto. Un módulo reutilizable no es un estado de la especificación, es
un tipo de especificación; borrador y validada, en cambio, sí son estados.

\subsec{06. Ampliación en detalle del proceso principal (SD1)}

El diagrama SD1 amplía el proceso principal en siete subprocesos. La
convención de la norma es que la posición vertical, y sólo ella, denota
precedencia temporal: no hay flechas entre subprocesos. Los objetos que ya
aparecen en el diagrama de sistema permanecen fuera de la elipse ampliada, y
los enlaces cruzan su contorno hasta el subproceso concreto que los usa, que
es lo que permite ver en qué punto exacto interviene cada uno. Los estados de
la \emph{Especificación de Infraestructura}, del \emph{Plan de Cambios} y de
la \emph{Infraestructura Desplegada} se hacen explícitos aquí, y con ellos las
transiciones.

\begin{figure}[H]
  \centering
  \begin{tikzpicture}[scale=0.79, transform shape, x=1mm, y=1mm,
                      every node/.style={transform shape}]
    \tikzset{sp/.style={opmproc, minimum width=46mm, minimum height=10mm,
               text width=, font=\footnotesize\bfseries},
             ob/.style={opmobj, text width=23mm, minimum height=8mm,
               font=\scriptsize\bfseries},
             obe/.style={draw=opmobjc, very thick, fill=opmobjc!7,
               minimum width=34mm},
             est/.style={opmstate, font=\tiny, minimum width=20mm}}

    % ---- elipse del proceso ampliado
    \node[draw=opmprocc, very thick, ellipse, fill=opmprocc!5,
      minimum width=100mm, minimum height=124mm] (marco) at (0,0) {};
    \node[font=\scriptsize\bfseries, align=center] at (0,50)
      {Aprovisionamiento\\de Infraestructura};

    % ---- los siete subprocesos: la posición vertical fija la precedencia
    \node[sp] (c)  at (0, 36) {Codificación};
    \node[sp] (v)  at (0, 24) {Validación};
    \node[sp] (pl) at (0, 12) {Planificación};
    \node[sp] (au) at (0,  0) {Autorización};
    \node[sp] (ap) at (0,-12) {Aplicación};
    \node[sp] (co) at (0,-24) {Conciliación};
    \node[sp] (re) at (0,-36) {Retirada};

    % ---- eje temporal, fuera de todo
    \draw[flecha, gray!65] (-106,42) -- (-106,-42)
      node[midway, rotate=90, above=0.6mm, font=\scriptsize,
           text=gray!55!black] {precedencia temporal};

    % ---- agente, fuera del proceso ampliado
    \node[ob, opmamb, opmfis, text width=21mm] (eq) at (-72,48)
      {Equipo de\\Plataforma};

    % ---- objetos de la izquierda
    \node[obe, minimum height=28mm] (spec) at (-72,28) {};
    \node[anchor=north, font=\scriptsize\bfseries, align=center]
      at ([yshift=-1.1mm]spec.north) {Especificación de\\Infraestructura};
    \node[est] (sb) at (-72,30) {borrador};
    \node[est] (sv) at (-72,22) {validada};

    \node[ob] (edo)   at (-72, 8) {Registro\\de Estado};
    \node[ob] (motor) at (-72,-6) {Motor de IaC};

    \node[obe, minimum height=40mm, opmfis] (infra) at (-72,-33) {};
    \node[anchor=north, font=\scriptsize\bfseries, align=center]
      at ([yshift=-1.1mm]infra.north) {Infraestructura\\Desplegada};
    \node[est] (si) at (-72,-24) {inexistente};
    \node[est] (sc) at (-72,-32) {conforme};
    \node[est] (sd) at (-72,-40) {desviada};
    \node[est] (sr) at (-72,-48) {retirada};

    % ---- objetos de la derecha
    \node[ob] (pol) at (72,26) {Conjunto de\\Políticas};
    \node[obe, minimum height=20mm] (plan) at (72,6) {};
    \node[anchor=north, font=\scriptsize\bfseries] at ([yshift=-1.1mm]plan.north)
      {Plan de Cambios};
    \node[est] (pp) at (72,7.5) {pendiente};
    \node[est] (pa) at (72,0.5) {aprobado};
    \node[ob] (cred) at (72,-16) {Credenciales};
    \node[ob, opmamb, opmfis] (prov) at (72,-30) {Proveedor\\de Nube};
    \node[ob] (aud)  at (72,-46) {Registro de\\Auditoría};

    % ---- enlaces del agente
    \draw[agente] (eq.south east) to[out=-40,in=155] (c.north west);
    \draw[agente, rounded corners=3mm]
      (eq.west) -- (-96,48) -- (-96,0) -- (au.west);

    % ---- codificación, validación y planificación
    \draw[flecha] (c.west) to[out=180,in=10] ([yshift=1mm]sb.east);
    \draw[flecha] ([yshift=-1mm]sb.east) to[out=-10,in=170] (v.north west);
    \draw[flecha] (v.west) to[out=180,in=10] ([yshift=1mm]sv.east);
    \draw[instr]  ([yshift=-1mm]sv.east) to[out=-10,in=165] (pl.north west);
    \draw[instr]  (edo.east) -- (pl.west);
    \draw[instr]  (motor.north east) to[out=25,in=195] (pl.south west);

    % ---- autorización y aplicación
    \draw[flecha] (pl.east) to[out=0,in=170] (pp.north west);
    \draw[instr]  (pp.south west) to[out=-165,in=15] (au.east);
    \draw[flecha] (au.north east) to[out=15,in=190] (pa.west);
    \draw[instr]  (pa.south west) to[out=-160,in=20] (ap.east);
    \draw[instr]  (motor.east) to[out=0,in=170] (ap.north west);
    \draw[instr]  (cred.west) -- (ap.east);
    \draw[instr]  (prov.north west) to[out=155,in=-15] (ap.south east);
    \draw[efecto] (ap.west) to[out=170,in=-25] (edo.south east);
    \draw[flecha] (si.east) to[out=0,in=200] (ap.south west);
    \draw[flecha] (ap.west) to[out=195,in=15] ([yshift=1.4mm]sc.east);

    % ---- conciliación y retirada
    \draw[instr]  (prov.west) to[out=180,in=-15] (co.east);
    \draw[flecha] (sd.east) to[out=0,in=195] (co.west);
    \draw[flecha] (co.north west) to[out=160,in=-5] ([yshift=-1.4mm]sc.east);
    \draw[flecha] (co.south east) to[out=-40,in=165] (aud.west);
    \draw[instr]  (prov.south west) to[out=205,in=-10] (re.east);
    \draw[flecha, rounded corners=2mm]
      ([yshift=-3.2mm]sc.east) -- (-49,-35.2) -- (-32,-35.2) -- (re.west);
    \draw[flecha] (re.south west) to[out=210,in=-10] (sr.south east);
  \end{tikzpicture}
  \caption{SD1: ampliación en detalle de \emph{Aprovisionamiento de
  Infraestructura} en sus siete subprocesos. La precedencia temporal la fija
  la posición vertical, sin flechas entre procesos. Los objetos que ya
  aparecen en el diagrama de sistema quedan fuera de la elipse y se enlazan
  con el subproceso que los usa. \emph{Aplicación}, \emph{Conciliación} y
  \emph{Retirada} actúan sobre el mismo objeto desde estados de partida
  distintos, que es la formulación precisa de la diferencia entre desplegar,
  corregir la deriva y desmantelar. Elaboración propia.}
  \label{fig:opmsd1}
\end{figure}

\subsec{07. El OPL de la ampliación}

\begin{quote}\small
\textsc{Aprovisionamiento de Infraestructura} se amplía en
\textsc{Codificación}, \textsc{Validación}, \textsc{Planificación},
\textsc{Autorización}, \textsc{Aplicación}, \textsc{Conciliación} y
\textsc{Retirada}, en ese orden temporal.\\
\textsc{Especificación de Infraestructura} puede estar borrador o validada.\\
\textsc{Plan de Cambios} puede estar pendiente o aprobado.\\
\textsc{Equipo de Plataforma} maneja \textsc{Codificación} y
\textsc{Autorización}.\\
\textsc{Codificación} produce \textsc{Especificación de Infraestructura}
borrador.\\
\textsc{Validación} requiere \textsc{Conjunto de Políticas}.\\
\textsc{Validación} cambia \textsc{Especificación de Infraestructura} de
borrador a validada.\\
\textsc{Planificación} requiere \textsc{Motor de IaC}, \textsc{Registro de
Estado} y \textsc{Especificación de Infraestructura} validada.\\
\textsc{Planificación} produce \textsc{Plan de Cambios} pendiente.\\
\textsc{Autorización} requiere \textsc{Plan de Cambios} pendiente.\\
\textsc{Autorización} cambia \textsc{Plan de Cambios} de pendiente a
aprobado.\\
\textsc{Aplicación} requiere \textsc{Motor de IaC}, \textsc{Plan de Cambios}
aprobado, \textsc{Credenciales} y \textsc{Proveedor de Nube}.\\
\textsc{Aplicación} cambia \textsc{Infraestructura Desplegada} de inexistente
a conforme.\\
\textsc{Aplicación} afecta a \textsc{Registro de Estado}.\\
\textsc{Conciliación} requiere \textsc{Proveedor de Nube}.\\
\textsc{Conciliación} cambia \textsc{Infraestructura Desplegada} de desviada a
conforme.\\
\textsc{Conciliación} produce \textsc{Registro de Auditoría}.\\
\textsc{Retirada} requiere \textsc{Proveedor de Nube}.\\
\textsc{Retirada} cambia \textsc{Infraestructura Desplegada} de conforme a
retirada.
\end{quote}

\subsec{08. Qué hace explícito el modelo}

Merece la pena señalar cuatro cosas que el modelo obliga a decir y que las
descripciones en prosa dejan implícitas.

\begin{enumerate}
  \item \textsc{Registro de Estado} es el único objeto que aparece a la vez
    como instrumento de la planificación y como objeto afectado por la
    aplicación. Esa doble condición es exactamente la razón por la que
    necesita bloqueo de concurrencia, y el modelo la hace visible sin
    argumentarla.
  \item \textsc{Aplicación}, \textsc{Conciliación} y \textsc{Retirada} cambian
    el mismo objeto desde estados de partida distintos. Con enlaces de efecto
    al objeto entero las tres serían indistinguibles; con enlaces de entrada y
    salida entre estados quedan formalmente separadas, que es la diferencia
    entre desplegar, corregir la deriva y desmantelar.
  \item El equipo humano interviene sólo en dos de los siete subprocesos, la
    codificación y la autorización. Eso delimita el alcance real de la
    automatización y, de paso, señala dónde está el cuello de botella de
    personas.
  \item La \textsc{Validación} se sitúa entre la \textsc{Codificación} y la
    \textsc{Planificación}, y la \textsc{Autorización} entre la
    \textsc{Planificación} y la \textsc{Aplicación}. Esa partición es la que
    rompe el bloque acoplado L1--M4 que la
    Figura~\ref{fig:dsmarqord} detecta en la matriz de dependencias: el
    veredicto de la política deja de necesitar un motor que a su vez lo
    necesita a él, porque cada uno actúa sobre un subproceso distinto.
\end{enumerate}

% ==========================================================================
\section{Figuras de mérito del sistema}
\label{sec:fom}
% ==========================================================================

Un sistema de IaC es en sí mismo una tecnología, y como tal admite figuras de
mérito que permiten situarlo y seguir su evolución. El
Cuadro~\ref{tab:fom} propone ocho, con su unidad y su procedimiento de
medida. Las cuatro primeras coinciden con las métricas de rendimiento de
entrega ampliamente utilizadas \parencite{forsgren2018}; las cuatro últimas
son específicas de la infraestructura.

\begin{table}[ht]
  \centering
  \footnotesize
  \caption{Figuras de mérito propuestas para un sistema de IaC.}
  \label{tab:fom}
  \begin{tabularx}{\textwidth}{@{}>{\raggedright\arraybackslash}p{0.27\textwidth}
    >{\raggedright\arraybackslash}p{0.17\textwidth}
    >{\raggedright\arraybackslash}X@{}}
    \toprule
    \textbf{Figura de mérito} & \textbf{Unidad} & \textbf{Cómo se mide} \\
    \midrule
    Tiempo de ciclo del cambio & horas & Del registro del cambio en el
      repositorio a su aplicación efectiva en producción. \\
    Frecuencia de aplicación & cambios/semana & Recuento de aplicaciones
      correctas por entorno. \\
    Tasa de fallo del cambio & \% & Aplicaciones que exigen reversión o
      corrección inmediata sobre el total. \\
    Tiempo de restablecimiento & minutos & De la detección del fallo a la
      recuperación del servicio. \\
    Cobertura de codificación & \% & Recursos gestionados por el motor sobre
      el total de recursos existentes en el proveedor. \\
    Tasa de deriva & recursos desviados / 1000 · semana & Diferencias
      detectadas en las ejecuciones de sólo lectura. \\
    Tiempo de reconstrucción & horas & De un entorno vacío a un entorno
      equivalente al de producción, partiendo sólo del repositorio. \\
    Cumplimiento de política & \% & Planes que superan la evaluación de
      políticas sin excepción manual. \\
    \bottomrule
  \end{tabularx}
\end{table}

El tiempo de reconstrucción es la figura más reveladora de todas, porque
verifica de un solo golpe las cuatro propiedades exigidas en la
sección~\ref{sec:alcance}. Una
organización que no puede reconstruir un entorno completo desde el repositorio
no tiene un sistema de IaC completo, con independencia de cuántos ficheros
declarativos posea.

% ==========================================================================
\section{Posicionamiento frente a la competencia}
\label{sec:comp}
% ==========================================================================

Esta sección corresponde al elemento 6 del índice de una hoja de ruta
tecnológica: el posicionamiento de la empresa con respecto a la competencia
expresado sobre las figuras de mérito de la tecnología
\parencite{deweck2022}. El instrumento empleado es el gráfico vectorial, que
sitúa a cada jugador en un plano cuyos ejes miden la variación de valor
percibida por el cliente y la variación de coste soportada por el productor, y
que descompone el desplazamiento de cada uno en los vectores tecnológicos
concretos que lo producen. Todo el análisis se apoya en las ocho figuras de
mérito definidas en el Cuadro~\ref{tab:fom}.

\subsec{01. Qué se compara y desde qué punto de vista}

Un gráfico vectorial obliga a fijar tres cosas antes de dibujar la primera
flecha, y ninguna de ellas es evidente cuando la tecnología analizada es una
capacidad interna y no un producto de catálogo.

\begin{enumerate}
  \item \textbf{El productor} es la organización de plataforma: el equipo que
    construye y opera el sistema de infraestructura como código descrito en
    las secciones anteriores. El coste que soporta es el coste anual de
    propiedad de ese sistema, no el coste de la infraestructura desplegada
    con él.
  \item \textbf{El cliente} es el equipo de producto que consume la
    plataforma. El valor que percibe no es abstracto: son las cuatro figuras
    de mérito del Cuadro~\ref{tab:fom} que se manifiestan en su trabajo
    diario, a saber, el tiempo de ciclo del cambio, la cobertura de
    codificación, la tasa de fallo del cambio y el tiempo de reconstrucción.
    Las otras cuatro son internas del productor y no entran en el eje de
    valor.
  \item \textbf{El producto de referencia} es la plataforma tal como está
    hoy, en 2026: motor de aprovisionamiento comunitario, estado remoto,
    canalización de integración continua, sin política como código, sin
    catálogo certificado y sin interfaz de autoservicio. Ocupa el origen de
    coordenadas por construcción, y todo el gráfico mide desviaciones
    respecto a él.
\end{enumerate}

Los jugadores comparados no son las empresas proveedoras entre sí, sino las
posiciones que alcanza una organización adoptante al construir su plataforma
sobre cada una de las tres familias tecnológicas al horizonte de 2029, es
decir, hoy más tres años. Esta distinción importa: la pregunta que responde el
gráfico no es cuál herramienta es mejor en abstracto, sino dónde queda situada
la empresa según cuál adopte.

\subsec{02. Los tres competidores considerados}

El Cuadro~\ref{tab:competidores} caracteriza a los tres jugadores. Se han
elegido porque cubren las tres estrategias de producto que hoy se disputan la
capa L1 de la Figura~\ref{fig:pila}: el estándar de hecho multinube
(Terraform), la integración vertical de un proveedor único (Bicep) y la
apuesta por los lenguajes de propósito general (Pulumi).

\begin{table}[ht]
  \centering
  \footnotesize
  \setlength{\tabcolsep}{4pt}
  \caption{Caracterización de los tres jugadores en la capa de
  aprovisionamiento.}
  \label{tab:competidores}
  \begin{tabularx}{\textwidth}{@{}>{\raggedright\arraybackslash}p{0.15\textwidth}
    >{\raggedright\arraybackslash}X
    >{\raggedright\arraybackslash}X
    >{\raggedright\arraybackslash}X@{}}
    \toprule
    & \textbf{Jugador A: Terraform} & \textbf{Jugador B: Bicep} &
      \textbf{Jugador C: Pulumi} \\
    \midrule
    Titular & HashiCorp, integrada en IBM desde 2025 & Microsoft &
      Pulumi Corporation \\
    Lenguaje & HCL, lenguaje declarativo propio
      \parencite{terraformdocs} & DSL declarativo que se
      traduce a plantillas ARM & lenguajes de propósito general: TypeScript,
      Python, Go, C\# \\
    Ámbito & multinube, con un registro de proveedores muy extenso &
      exclusivamente Azure & multinube, reutiliza los proveedores de
      Terraform \\
    Estado & fichero de estado explícito, remoto, cifrado y con bloqueo
      \parencite{terraformstate} & sin fichero de estado: lo mantiene el
      gestor de recursos de Azure \parencite{deploymentstacks} & fichero de
      estado, en servicio gestionado o autoalojado \\
    Licencia & BUSL desde 2023 \parencite{hashicorpbusl}; servicio gestionado
      de pago & gratuito, incluido en la suscripción de Azure
      \parencite{bicepdocs} & núcleo abierto; servicio gestionado con tarifa
      por recurso \parencite{pulumidocs} \\
    Política como código & Sentinel u OPA \parencite{sentinel,opa} &
      Azure Policy, integrado en el plano de control &
      CrossGuard \\
    Reutilización & registro de módulos, público y privado & módulos
      verificados publicados por el propio proveedor \parencite{avm} &
      componentes que son clases del lenguaje anfitrión \\
    Barrera de adopción & expresividad limitada de HCL y coste de licencia al
      crecer & dependencia de un proveedor único & exige competencia real de
      programación en el equipo de plataforma \\
    \bottomrule
  \end{tabularx}
\end{table}

Junto a los tres jugadores se sitúa la ruta propuesta para la empresa, que no
coincide con ninguno de ellos porque combina un motor compatible de código
abierto \parencite{opentofu} con dos inversiones propias, el catálogo de
módulos certificados y la interfaz de autoservicio. El gráfico incluye además
una quinta posición, la contención de coste, que se evaluó y se descartó, y
cuya única función aquí es mostrar que el cuadrante inferior izquierdo del
plano existe y es alcanzable.

\subsec{03. Construcción de las dos coordenadas}

Los dos ejes del gráfico no son magnitudes observables sino índices
construidos, y conviene declarar cómo se calculan antes de leer ninguna
posición.

\textbf{Eje de ordenadas: variación de valor para el cliente.} Cada una de las
cuatro figuras de mérito orientadas al cliente se convierte en una utilidad
acotada entre 0 y 100 mediante interpolación lineal entre dos anclas, el peor
valor aceptable y el mejor valor alcanzable con la tecnología disponible, al
modo de la exploración de espacios de diseño multiatributo
\parencite{ross2004}. Para las figuras en las que menos es mejor,

\[
u_i = 100 \cdot \frac{x_i^{\text{peor}} - x_i}
                     {x_i^{\text{peor}} - x_i^{\text{mejor}}},
\]

y con el numerador invertido para la cobertura de codificación, que es la
única en la que más es mejor. Las anclas adoptadas son: tiempo de ciclo, de
48 a 4 horas; cobertura, del 40 al 95 \%; tasa de fallo, del 20 al 3 \%;
tiempo de reconstrucción, de 72 a 6 horas. El valor agregado es la media
ponderada

\[
V = 0{,}30\,u_{\text{ciclo}} + 0{,}25\,u_{\text{cobertura}}
  + 0{,}25\,u_{\text{fallo}} + 0{,}20\,u_{\text{reconstrucción}},
\]

y la coordenada del gráfico es $\Delta V = V - V_{\text{ref}}$, con
$V_{\text{ref}} = 40{,}3$ puntos para el producto de referencia. Los pesos
reflejan una decisión explícita: se prima la velocidad percibida sobre la
capacidad de reconstrucción, porque la primera se experimenta cada día y la
segunda sólo ante un incidente grave.

\textbf{Eje de abscisas: variación de coste para el productor.} Es un índice
de coste anual de propiedad por cada cien recursos gestionados, con la
plataforma actual normalizada a 100 puntos. Se descompone en tres sumandos:
licencia y servicio gestionado, ingeniería de mantenimiento del propio sistema
de plataforma, y formación e incorporación de personas. La coordenada del
gráfico es $\Delta C = C - 100$.

Una advertencia sobre el eje de coste: mide el régimen permanente en 2029, no
la inversión de transición. El coste de una sola vez de cada migración
pertenece al modelo financiero, que es el elemento 8 del índice y queda fuera
de esta sección. Un jugador puede parecer barato en el gráfico y exigir un
desembolso inicial considerable para llegar a esa posición.

\subsec{04. Estimación de las figuras de mérito al horizonte 2029}

El Cuadro~\ref{tab:fomcomp} recoge las estimaciones y las utilidades
resultantes. Son estimaciones razonadas de ingeniería, no medidas: se apoyan
en el comportamiento observado de la plataforma actual y en las propiedades
declaradas de cada tecnología, y su función es ordenar posiciones, no predecir
cifras.

\begin{table}[ht]
  \centering
  \footnotesize
  \setlength{\tabcolsep}{4.5pt}
  \caption{Figuras de mérito estimadas al horizonte 2029 y valor agregado
  resultante. Las cuatro figuras son las orientadas al cliente del
  Cuadro~\ref{tab:fom}.}
  \label{tab:fomcomp}
  \begin{tabular}{@{}lrrrrrr@{}}
    \toprule
    \textbf{Posición} & \textbf{Ciclo} & \textbf{Cobert.} &
      \textbf{Fallo} & \textbf{Reconstr.} & \textbf{$V$} &
      \textbf{$\Delta V$} \\
    & (h) & (\pct) & (\pct) & (h) & & \\
    \midrule
    Producto de referencia (2026) & 26 & 55 & 14 & 40 & 40,3 & 0,0 \\
    \midrule
    A. Terraform gestionado       &  9 & 85 &  6 & 12 & 85,8 & $+45{,}5$ \\
    B. Bicep sobre Azure          & 12 & 72 &  8 & 16 & 73,7 & $+33{,}4$ \\
    C. Pulumi                     &  7 & 88 &  5 & 10 & 90,6 & $+50{,}3$ \\
    \textbf{Ruta propia}          &  8 & 86 &  6 & 11 & 87,3 & $+46{,}9$ \\
    \midrule
    D. Contención de coste (descartada) & 32 & 44 & 15 & 60 & 23,7 &
      $-16{,}6$ \\
    \bottomrule
  \end{tabular}
\end{table}

El Cuadro~\ref{tab:coste} descompone el otro eje. La diferencia decisiva no
está en la licencia sino en la ingeniería de mantenimiento: un catálogo de
módulos certificados y una política como código reducen el trabajo recurrente
de revisión y corrección lo bastante para compensar el esfuerzo de
construirlos y sostenerlos.

\begin{table}[ht]
  \centering
  \footnotesize
  \setlength{\tabcolsep}{5pt}
  \caption{Descomposición del índice de coste anual de propiedad por cada cien
  recursos gestionados. La plataforma actual se normaliza a 100 puntos.}
  \label{tab:coste}
  \begin{tabular}{@{}lrrrrr@{}}
    \toprule
    \textbf{Posición} & \textbf{Licencia y} & \textbf{Ingeniería de} &
      \textbf{Formación} & \textbf{Total} & \textbf{$\Delta C$} \\
    & \textbf{servicio} & \textbf{mantenimiento} & & & \\
    \midrule
    Producto de referencia (2026) &  8 & 76 & 16 & 100 & 0 \\
    \midrule
    A. Terraform gestionado       & 52 & 60 & 20 & 132 & $+32$ \\
    B. Bicep sobre Azure          &  0 & 54 & 22 &  76 & $-24$ \\
    C. Pulumi                     & 46 & 62 & 36 & 144 & $+44$ \\
    \textbf{Ruta propia}          & 12 & 66 & 18 &  96 & $-4$ \\
    \midrule
    D. Contención de coste (descartada) & 6 & 56 & 12 & 74 & $-26$ \\
    \bottomrule
  \end{tabular}
\end{table}

La Figura~\ref{fig:perfilfom} compara los perfiles de utilidad figura por
figura, antes de agregarlos. Muestra algo que la coordenada agregada esconde:
las cuatro posiciones avanzadas son casi indistinguibles en velocidad y en
fiabilidad, y toda su diferencia se concentra en la cobertura de codificación.
Bicep no pierde por ser lento, sino por dejar recursos fuera del alcance del
motor.

\begin{figure}[ht]
  \centering
  \begin{tikzpicture}
    \begin{axis}[
      width=15.2cm, height=6.4cm,
      ybar=1pt, bar width=7pt,
      ymin=0, ymax=104,
      ylabel={\footnotesize Utilidad (0 a 100)},
      ylabel style={font=\footnotesize},
      symbolic x coords={f1,f2,f3,f4},
      xtick={f1,f2,f3,f4},
      xticklabels={Tiempo de ciclo, Cobertura de codificación,
                   Tasa de fallo, Tiempo de reconstrucción},
      xticklabel style={font=\scriptsize},
      yticklabel style={font=\scriptsize},
      ytick={0,20,40,60,80,100},
      ymajorgrids, grid style={gray!22},
      axis lines=left,
      enlarge x limits=0.16,
      legend style={font=\scriptsize, at={(0.5,1.04)}, anchor=south,
                    legend columns=5, draw=none, column sep=1.1ex},
    ]
      \addplot[fill=gray!40,  draw=gray!70!black]
        coordinates {(f1,50.0) (f2,27.3) (f3,35.3) (f4,48.5)};
      \addplot[fill=jugA!35,  draw=jugA]
        coordinates {(f1,88.6) (f2,81.8) (f3,82.4) (f4,90.9)};
      \addplot[fill=jugB!35,  draw=jugB]
        coordinates {(f1,81.8) (f2,58.2) (f3,70.6) (f4,84.8)};
      \addplot[fill=jugC!30,  draw=jugC]
        coordinates {(f1,93.2) (f2,87.3) (f3,88.2) (f4,93.9)};
      \addplot[fill=propia!35, draw=propia]
        coordinates {(f1,90.9) (f2,83.6) (f3,82.4) (f4,92.4)};
      \legend{Referencia 2026, Terraform, Bicep, Pulumi, Ruta propia}
    \end{axis}
  \end{tikzpicture}
  \caption{Perfil de utilidad de las cuatro figuras de mérito orientadas al
  cliente, antes de la agregación ponderada. La única diferencia relevante
  entre las cuatro posiciones avanzadas está en la cobertura de codificación.
  Elaboración propia.}
  \label{fig:perfilfom}
\end{figure}

\subsec{05. Descomposición en vectores tecnológicos}

El gráfico vectorial no dibuja las posiciones finales como puntos aislados
sino como la suma de los saltos tecnológicos que conducen a ellas. El
Cuadro~\ref{tab:vectores} identifica cada vector, la inversión que representa
y su contribución a las dos coordenadas.

\begin{table}[ht]
  \centering
  \footnotesize
  \setlength{\tabcolsep}{5pt}
  \caption{Descomposición de cada trayectoria en vectores tecnológicos. Las
  contribuciones son aditivas y suman la posición final de cada jugador.}
  \label{tab:vectores}
  \begin{tabularx}{\textwidth}{@{}l>{\raggedright\arraybackslash}X rr@{}}
    \toprule
    \textbf{Vector} & \textbf{Inversión tecnológica que representa} &
      \textbf{$\Delta C$} & \textbf{$\Delta V$} \\
    \midrule
    \multicolumn{4}{@{}l}{\textit{Jugador A: Terraform}} \\
    A1 & Ejecución y estado gestionados en el servicio del proveedor, con
         control de acceso por espacio de trabajo & $+24$ & $+30{,}0$ \\
    A2 & Despliegue multientorno por pilas y política como código con Sentinel
       & $+8$ & $+15{,}5$ \\
    \midrule
    \multicolumn{4}{@{}l}{\textit{Jugador B: Bicep}} \\
    B1 & Sustitución de las plantillas ARM por el DSL nativo, que suprime el
         fichero de estado y su gobierno & $-20$ & $+14{,}0$ \\
    B2 & Pilas de despliegue, módulos verificados del proveedor y Azure Policy
         en el plano de control & $-4$ & $+19{,}4$ \\
    \midrule
    \multicolumn{4}{@{}l}{\textit{Jugador C: Pulumi}} \\
    C1 & Descripción de la infraestructura en un lenguaje de propósito general
         con pruebas unitarias reales & $+32$ & $+26{,}0$ \\
    C2 & Gestión de entornos y secretos y interfaz programable de automatización
         \parencite{pulumiesc} & $+12$ & $+24{,}3$ \\
    \midrule
    \multicolumn{4}{@{}l}{\textit{Ruta propia}} \\
    P1 & Catálogo de módulos certificados y política como código con OPA
         \parencite{opa} & $+6$ & $+21{,}0$ \\
    P2 & Sustitución del motor por la bifurcación abierta y estado gestionado
         propio \parencite{opentofu} & $-18$ & $+1{,}9$ \\
    P3 & Interfaz de autoservicio construida sobre el catálogo
         \parencite{backstage} & $+8$ & $+24{,}0$ \\
    \midrule
    \multicolumn{4}{@{}l}{\textit{Opción descartada}} \\
    D1 & Contención de coste: reducir el alcance codificado al núcleo crítico
         y devolver el resto a operación manual & $-26$ & $-16{,}6$ \\
    \bottomrule
  \end{tabularx}
\end{table}

El orden de los vectores de la ruta propia no es arbitrario, y es la decisión
más discutible de toda la sección. El catálogo de módulos se construye
\emph{antes} de sustituir el motor, y no después, porque el catálogo es
precisamente lo que abarata la sustitución: cuando los equipos de producto
consumen módulos y no recursos, el motor queda detrás de una interfaz estable
y cambiarlo deja de ser un proyecto de migración para convertirse en un
cambio interno del catálogo. Invertir el orden convertiría P2 en el vector más
caro y arriesgado de los tres.

\subsec{06. El gráfico vectorial}

La Figura~\ref{fig:vectorial} compone los vectores del
Cuadro~\ref{tab:vectores} sobre el plano de las dos coordenadas. Se lee así:
todas las trayectorias arrancan del origen, que es la plataforma actual; cada
flecha es una inversión tecnológica concreta y se encadena a la anterior, de
modo que la posición de un jugador es la suma vectorial de sus inversiones y
no un punto postulado; el punto grueso del final de cada cadena es la posición
en 2029; y la flecha discontinua es la opción que se evaluó y se descartó.
Cuanto más arriba, más valor recibe el equipo de producto; cuanto más a la
izquierda, menos cuesta sostener la plataforma. El cuadrante superior
izquierdo es, por tanto, el favorable.

\begin{figure}[ht]
  \centering
  \begin{tikzpicture}[x=0.125cm, y=0.113cm,
      vec/.style  = {line width=1.05pt, -{Stealth[length=5.5pt]}},
      pto/.style  = {circle, fill, inner sep=2.1pt},
      cod/.style  = {font=\scriptsize\bfseries, inner sep=1.6pt},
      jugl/.style = {font=\footnotesize\bfseries, align=center, inner sep=2pt}]

    % rejilla de diez puntos de índice
    \foreach \x in {-40,-30,-20,-10,10,20,30,40,50}
      \draw[gray!20, thin] (\x,-24) -- (\x,56);
    \foreach \y in {-20,-10,10,20,30,40,50}
      \draw[gray!20, thin] (-42,\y) -- (58,\y);

    % ejes
    \draw[thick,-{Stealth[length=6pt]}] (-44,0) -- (62,0);
    \draw[thick,-{Stealth[length=6pt]}] (0,-26) -- (0,60);
    \node[font=\footnotesize\bfseries, align=center, anchor=south]
      at (2,61) {$\Delta$ Valor\\para el cliente};
    \node[font=\footnotesize\bfseries, align=center, anchor=north]
      at (52,-5) {$\Delta$ Coste para\\el productor};

    % marcas de escala
    \foreach \x in {-40,-30,-20,20,30,40,50}
      \node[font=\tiny, text=gray!55!black, below=1pt] at (\x,0) {\x};
    \foreach \y in {-10,10,20,30,40}
      \node[font=\tiny, text=gray!55!black, left=1pt] at (0,\y) {\y};

    \node[font=\scriptsize\itshape, text=gray!60!black, anchor=north west]
      at (-43,60) {Referencia: 2026 + 3 años};

    % ------------------------------------------------ jugador A: Terraform
    \draw[vec, jugA] (0,0) -- (24,30)
      node[cod, text=jugA, pos=0.55, below right=-1pt] {A1};
    \draw[vec, jugA] (24,30) -- (32,45.5)
      node[cod, text=jugA, midway, right=0pt] {A2};
    \node[pto, jugA] at (32,45.5) {};
    \node[jugl, text=jugA, anchor=east] at (29,45.5) {Jugador A\\Terraform};

    % ------------------------------------------------ jugador C: Pulumi
    \draw[vec, jugC] (0,0) -- (32,26)
      node[cod, text=jugC, pos=0.6, below right=-1pt] {C1};
    \draw[vec, jugC] (32,26) -- (44,50.3)
      node[cod, text=jugC, pos=0.32, right=0pt] {C2};
    \node[pto, jugC] at (44,50.3) {};
    \node[jugl, text=jugC, anchor=south] at (44,53) {Jugador C\\Pulumi};

    % ------------------------------------------------ jugador B: Bicep
    \draw[vec, jugB] (0,0) -- (-20,14)
      node[cod, text=jugB, pos=0.55, above right=-2pt] {B1};
    \draw[vec, jugB] (-20,14) -- (-24,33.4)
      node[cod, text=jugB, midway, right=0pt] {B2};
    \node[pto, jugB] at (-24,33.4) {};
    \node[jugl, text=jugB, anchor=east] at (-27,33.4) {Jugador B\\Bicep};

    % ------------------------------------------------ ruta propia
    \draw[vec, propia] (0,0) -- (6,21)
      node[cod, text=propia, pos=0.6, right=0pt] {P1};
    \draw[vec, propia] (6,21) -- (-12,22.9)
      node[cod, text=propia, pos=0.78, below=1pt] {P2};
    \draw[vec, propia] (-12,22.9) -- (-4,46.9)
      node[cod, text=propia, midway, left=0pt] {P3};
    \node[pto, propia] at (-4,46.9) {};
    \node[jugl, text=propia, anchor=east] at (-7,50)
      {Nuestra empresa\\ruta propuesta};

    % ------------------------------------------------ opción descartada
    \draw[vec, descart, dashed] (0,0) -- (-26,-16.6)
      node[cod, text=descart, pos=0.5, below=1pt] {D1};
    \node[pto, descart] at (-26,-16.6) {};
    \node[font=\scriptsize\itshape, text=descart, align=center, anchor=north]
      at (-26,-19.5) {Contención de coste\\(descartada)};

    % ------------------------------------------------ origen
    \node[circle, draw, very thick, fill=white, inner sep=1.7pt] at (0,0) {};
    \node[font=\scriptsize, align=left, anchor=north west] at (7,-7)
      {Producto de referencia:\\plataforma actual, 2026};
  \end{tikzpicture}
  \caption{Gráfico vectorial de posicionamiento competitivo al horizonte 2029.
  El origen es la plataforma actual. Cada flecha es un vector tecnológico del
  Cuadro~\ref{tab:vectores} y cada punto, la posición que alcanza una
  organización que adopta esa trayectoria. La rejilla es de diez puntos de
  índice. Elaboración propia.}
  \label{fig:vectorial}
\end{figure}

\subsec{07. Lectura del gráfico}

El gráfico dice cuatro cosas, y ninguna de ellas es que exista una herramienta
mejor que las demás.

\textbf{Primera: la frontera de eficiencia la forman tres posiciones, no una.}
Entre las cuatro trayectorias de inversión, Pulumi, la ruta propia y Bicep no
se dominan entre sí. Pulumi ofrece el mayor
valor, $+50{,}3$ puntos, y cuesta $+44$; Bicep es el más barato, $-24$, y
entrega $+33{,}4$; la ruta propia se sitúa en $-4$ y $+46{,}9$. Terraform
gestionado, en cambio, sí está dominado: la ruta propia entrega más valor
($+46{,}9$ frente a $+45{,}5$) por 36 puntos menos de coste. Ésa es la
conclusión operativa más limpia de todo el análisis, y explica por qué el
vector P2 existe.

\textbf{Segunda: el valor no lo aporta el motor, lo aportan las capas de
encima.} En la ruta propia, el vector que sustituye el motor, P2, aporta
$+1{,}9$ puntos de valor sobre 46,9, es decir, un 4 \%. Los otros 45 puntos
vienen del catálogo y del autoservicio, que son inversiones propias y no
dependen del proveedor elegido. La consecuencia práctica es que la decisión de
herramienta, que suele consumir la mayor parte del debate técnico, es la menos
determinante de las tres.

\textbf{Tercera: la Figura~\ref{fig:perfilfom} explica de dónde sale la
distancia.} Las cuatro posiciones avanzadas están entre 81 y 93 puntos de
utilidad en velocidad y entre 70 y 88 en fiabilidad, márgenes estrechos. La
cobertura de codificación, en cambio, va de 58 a 87. Lo que separa a los
jugadores no es lo bien que hacen aquello que hacen, sino cuánto de la
infraestructura queda dentro de su alcance.

\textbf{Cuarta: el cuadrante inferior izquierdo es una trampa reconocible.} La
opción descartada D1 ahorra 26 puntos de coste, dos puntos más que Bicep, y
destruye 16,6 puntos de valor. Conviene ser exacto aquí: al ser la más barata
de todas, D1 tampoco está dominada, y una regla de eficiencia de Pareto
aplicada sin más la conservaría en la frontera. Se descarta por un criterio
que el plano no contiene, un suelo de valor aceptable, y esto ilustra el
límite del instrumento: el gráfico ordena posiciones, pero no decide cuál es
admisible. Reducir el alcance codificado no elimina el
trabajo de infraestructura, lo devuelve a la forma manual, donde no se mide y
no aparece en el índice de coste de la plataforma porque se ha desplazado a
los equipos de producto. El vector es barato en el eje que el gráfico mide y caro en uno
que no mide.

\subsec{08. Sensibilidad al supuesto de cartera}

Toda la comparación descansa en un supuesto que conviene poner a prueba: que
la cartera de infraestructura seguirá concentrada en un solo proveedor de
nube. El Cuadro~\ref{tab:sensib} recalcula las cuatro posiciones bajo el
escenario contrario, con un 40 \% de la carga fuera de Azure, que obliga al
jugador B a sostener una segunda cadena de herramientas.

\begin{table}[ht]
  \centering
  \footnotesize
  \setlength{\tabcolsep}{6pt}
  \caption{Efecto del escenario multinube sobre las posiciones. Sólo el
  jugador B se desplaza de forma apreciable.}
  \label{tab:sensib}
  \begin{tabular}{@{}lrrrrrr@{}}
    \toprule
    & \multicolumn{2}{c}{\textbf{Proveedor único}} &
      \multicolumn{2}{c}{\textbf{Cartera multinube}} &
      \multicolumn{2}{c}{\textbf{Desplazamiento}} \\
    \cmidrule(lr){2-3}\cmidrule(lr){4-5}\cmidrule(lr){6-7}
    \textbf{Posición} & $\Delta C$ & $\Delta V$ & $\Delta C$ & $\Delta V$ &
      $\Delta C$ & $\Delta V$ \\
    \midrule
    A. Terraform gestionado & $+32$ & $+45{,}5$ & $+35$ & $+43{,}6$ &
      $+3$ & $-1{,}9$ \\
    B. Bicep sobre Azure    & $-24$ & $+33{,}4$ & $-4$  & $+10{,}9$ &
      $+20$ & $-22{,}5$ \\
    C. Pulumi               & $+44$ & $+50{,}3$ & $+47$ & $+48{,}4$ &
      $+3$ & $-1{,}9$ \\
    \textbf{Ruta propia}    & $-4$  & $+46{,}9$ & $-1$  & $+45{,}0$ &
      $+3$ & $-1{,}9$ \\
    \bottomrule
  \end{tabular}
\end{table}

El resultado es asimétrico y por eso es útil. Las tres posiciones multinube se
mueven tres puntos en coste y dos en valor, un desplazamiento despreciable. El
jugador B pierde 22,5 puntos de valor y sube 20 de coste, hasta dejar de ser
la opción barata: cae de la frontera de eficiencia al interior del plano,
dominado por la ruta propia en ambos ejes. Su ventaja competitiva no es una
propiedad de la tecnología sino una propiedad del supuesto de cartera, y
desaparece con él.

Esto convierte la elección en una apuesta explícita y no en una preferencia
técnica: adoptar Bicep es afirmar que la empresa no abrirá una segunda nube en
el horizonte del plan. Si esa afirmación no puede sostenerse por escrito ante
la dirección, el jugador B queda fuera con independencia de lo que diga el
gráfico principal.

\subsec{09. Límites del análisis}

Tres reservas que conviene dejar escritas junto al gráfico, porque un plano
con flechas transmite más certidumbre de la que el método admite.

\begin{enumerate}
  \item \textbf{Las coordenadas son estimaciones, no medidas.} Las figuras de
    mérito del Cuadro~\ref{tab:fomcomp} proceden del comportamiento observado
    de la plataforma actual y de las propiedades declaradas de cada
    tecnología. El propio Cuadro~\ref{tab:fom} define cómo medirlas: la
    manera correcta de validar este análisis es instrumentar la referencia y
    volver a dibujar el gráfico con datos de un trimestre real.
  \item \textbf{Los pesos del índice de valor son una decisión de gobierno.}
    Al ponderar el tiempo de ciclo con 0,30 y el de reconstrucción con 0,20 se
    está escogiendo una postura ante el riesgo. Una organización sujeta a un
    requisito de recuperación ante desastres invertiría esos pesos, y con
    ellos el orden de los dos primeros clasificados.
  \item \textbf{El eje de coste ignora la transición.} Como se advirtió en el
    apartado 03, mide el régimen permanente de 2029. La ruta propia exige tres
    inversiones sucesivas cuyo desembolso inicial no aparece en ninguna
    coordenada, y esa cifra es la que decide si el plan es financiable. Su
    lugar es el elemento 8 del índice, el modelo financiero.
\end{enumerate}

% ==========================================================================
\section{Conclusiones}
\label{sec:conclusiones}
% ==========================================================================

\begin{enumerate}
  \item El funcionamiento de un sistema de infraestructura como código se
    reduce a un mecanismo único: la conciliación repetida entre un estado
    deseado, expresado en código versionado, y un estado real, observado a
    través de la interfaz de un proveedor. Todo lo demás, herramientas,
    formatos y flujos de trabajo, son variaciones sobre ese mecanismo.
  \item La pieza que hace posible el mecanismo es el estado registrado. No es
    un detalle de implementación sino el elemento que permite reconocer la
    propiedad de los recursos, detectar destrucciones y encadenar dependencias,
    y por eso concentra los requisitos de cifrado, versionado y bloqueo.
  \item Un sistema completo comprende siete capas, del sustrato del proveedor
    a la interfaz de autoservicio, y seis módulos transversales. Las capas
    intermedias compiten por el mismo territorio, de modo que la primera
    decisión de arquitectura consiste en asignar de forma explícita la
    propiedad de cada característica a una sola capa.
  \item Los módulos transversales, y en particular la política como código y
    las pruebas, son los que distinguen un sistema que automatiza de uno que
    además garantiza. Su ausencia no impide desplegar: hace que los errores se
    desplieguen igual de rápido que los aciertos.
  \item El tiempo de reconstrucción de un entorno completo desde el
    repositorio es la figura de mérito que mejor resume el estado de una
    implantación, porque verifica de una vez las cuatro propiedades que
    definen al sistema.
  \item La matriz de estructura de dependencia desplaza el orden de
    implantación que sugiere la intuición. Los concentradores de riesgo no son
    las capas visibles sino dos módulos transversales, el control de versiones
    con ocho salidas y la gestión de secretos con seis, y en la secuencia
    particionada la gestión de secretos aparece en la segunda ola, antes que
    cualquier capa intermedia. Tratarla como un endurecimiento posterior
    obliga a rehacer todo lo construido antes.
  \item La matriz de la cartera resulta enteramente triangular inferior, de
    modo que existe una secuencia de los ocho proyectos que no obliga a
    rehacer ninguno: custodia del estado, cierre de cobertura, catálogo y
    observabilidad, política y pruebas, sustitución del motor y, en último
    lugar, interfaz de autoservicio. Convierte además en constatación
    estructural lo que la sección~\ref{sec:comp} defiende como criterio: el
    catálogo antes que el motor no es una preferencia sino la única
    ordenación que no fuerza una iteración.
  \item El gráfico vectorial de la sección~\ref{sec:comp} desplaza la
    pregunta habitual.
    La comparación entre Terraform, Bicep y Pulumi ocupa la mayor parte del
    debate técnico y explica menos de la décima parte del valor que se puede
    ganar: en la trayectoria propuesta, el vector que sustituye el motor
    aporta 1,9 de los 46,9 puntos de mejora, y los otros 45 proceden del
    catálogo de módulos certificados y de la interfaz de autoservicio, que son
    inversiones propias e independientes del proveedor elegido.
  \item Lo que separa a los jugadores en el plano de las figuras de mérito no
    es la velocidad ni la fiabilidad, magnitudes en las que las cuatro
    posiciones avanzadas quedan a pocos puntos entre sí, sino la cobertura de
    codificación, esto es, cuánta infraestructura queda dentro del alcance del
    motor. La ventaja de coste del jugador B no es una propiedad de su
    tecnología sino del supuesto de proveedor único: bajo una cartera
    multinube pierde 22,5 puntos de valor y abandona la frontera de
    eficiencia.
  \item El modelado en OPM aporta algo que la descripción en prosa no aporta:
    obliga a declarar el tipo de cada relación. Al hacerlo aparecen con
    claridad la doble condición del registro de estado, la diferencia formal
    entre aplicar, conciliar y retirar, que sólo se sostiene si los enlaces
    parten de estados y no del objeto entero, y el alcance exacto de la
    intervención humana, que en este modelo se limita a dos de los siete
    subprocesos.
  \item Los elementos 2 y 3 del índice se sostienen mutuamente, y ese es el
    hallazgo metodológico del trabajo. La matriz de dependencias detecta un
    único bloque acoplado, el que forman el motor de aprovisionamiento y la
    política como código, y no dispone de instrumentos para romperlo; la
    ampliación en detalle del modelo OPM lo rompe al partir el proceso en
    planificación, validación y aplicación. La matriz encuentra el problema y
    el modelo aporta la partición que lo resuelve.
\end{enumerate}

\printbibliography[heading=bibintoc,title={Referencias bibliográficas}]

\end{document}
